% The WFDB library functions: the reference of the C-style API of WFDB-x.
% Part of the WFDB-x user's manual (see ../README.md for the conventions).
% check_manual.py checks that every library function of the Python
% implementation (the C-named functions and the wfdbx_* additions) has
% exactly one \apientry here, and that every entry is one of them.

\chapter{The WFDB library functions}
\label{chap:library}
% (long identifiers: allow more stretch in this chapter; reset at its end)
\setlength{\emergencystretch}{2.5em}

Besides its high-level functions, which read and write whole records in the
idiom of each language, WFDB-x provides the functions of the C library of the
WFDB Software Package \wfdbversion{} (\code{libwfdb}, declared in
\file{<wfdb/wfdb.h>}) in all seven languages. They keep the C names
(\code{isigopen}, \code{getvec}, \code{getann}, \code{strtim}, ...), the C
arguments and their order, the C return values and error codes, the messages
that the C library prints, and its global state: the open record, signals and
annotators, the sampling frequency and base time used by the time conversion
functions, the calibration list, the tables of annotation mnemonics. A program
written for the C library, or a passage of the \emph{WFDB Programmer's Guide},
carries over call by call; the command-line tools of WFDB-x
(\cref{chap:tools}) are written with these functions, as the C tools are with
\code{libwfdb}.

The behavior is that of the C library, edge cases included, with the
deliberate changes of WFDB-x described in \cref{chap:behavior}: no built-in
sampling frequency, gain, ADC resolution or units; files found only where
they are named (an input directory, an output directory, network access off
by default) instead of being searched for in the WFDB path; out-of-order
annotation files sorted in the same process, without losing annotations; an
output signal file never one of the input files; write errors counted. Legacy
mode (\code{wfdbx\_set\_legacy}) restores the original behavior. Each entry
below says what changes, if anything. WFDB-x adds a few functions of its own,
whose names begin with \code{wfdbx\_} (\cref{sec:library-additions}), and a
reader of many sample vectors at once (\code{getframes}, or \code{getvecs}).

Since the names, the arguments and the behavior are the same in the seven
languages, each function is documented once, in C notation: the first listing
of each entry is the prototype of the C library (for the WFDB-x additions,
the declaration of the C++ implementation, \file{cpp/include/wfdbx/wfdb.hpp},
which is in the same style). \Cref{sec:library-calling} shows how the C
conventions (arrays filled by a function, output arguments, \code{NULL},
strings, structures) appear in each language. The record \code{100s} used in
the examples is the sample record included with WFDB-x (the first minute of
record 100 of the MIT-BIH Arrhythmia Database, in \file{tests/data}).

\begin{note}
These functions are for programs that need the C library's model: reading
sample by sample, writing headers and annotation files piece by piece, or
porting an existing C program. To read or write a whole record or annotation
file, the high-level functions of each language (\code{read\_record},
\code{write\_annotations}, ...; \cref{chap:python,chap:cpp,chap:rust,chap:julia,chap:r,chap:matlab,chap:kotlin})
are simpler and faster.
\end{note}

% ==========================================================================
\section{Calling the library functions from each language}
\label{sec:library-calling}

In every language the functions have their C names and take their arguments
in the C order. What differs is how a language expresses C's pointers:
arrays that a function fills, output arguments, \code{NULL}, and strings.
\Cref{tab:library-languages} summarizes these conventions; the notes and the
example that follow give the details for each language. Each example opens
the record \code{100s} (from the current directory, the default input
directory) with \code{isigopen}, reads one sample of each signal with
\code{getvec}, and prints the description of signal 0 and the two samples
(\code{MLII 995 1011}).

The conventions common to all languages:

\begin{itemize}
\item Signal, group and annotator numbers start at 0, as in C, also in the
  languages whose arrays start at 1 (Julia, R, MATLAB/Octave).
\item Errors are reported as in C: by the return value, with a message on
  the standard error output, not by exceptions. \code{wfdberror} returns the
  last message, and \code{wfdbquiet} stops the messages from being printed.
\item The state of the library is kept between calls, as in C.
  \code{wfdbquit} closes the files and resets the record's state; the input
  and output directories and the network setting are kept.
\item Text is handled as bytes, one character per byte (ISO-8859-1 where a
  language's strings are Unicode), so that every byte of a header or an
  annotation file is kept.
\item An annotation's \code{aux} field is either absent (C's \code{NULL}) or a
  sequence of bytes whose first byte is the length of the data that follow,
  as in C.
\end{itemize}

\begin{xltabular}{\linewidth}{@{}P{1.6cm}LLLL@{}}
\caption{The C conventions in each language.}
\label{tab:library-languages}\\
\toprule
{\tabhead Language} & {\tabhead Calling} & {\tabhead Arrays and output
arguments} & {\tabhead \code{NULL}} & {\tabhead Strings returned}\\
\midrule
\endfirsthead
\toprule
{\tabhead Language} & {\tabhead Calling} & {\tabhead Arrays and output
arguments} & {\tabhead \code{NULL}} & {\tabhead Strings returned}\\
\midrule
\endhead
\bottomrule
\endfoot
Python & \code{wfdbx.isigopen(...)} (module \code{wfdbx}) & a list passed by
the caller is filled; an object (\code{Annotation}, \code{Calinfo}) passed by
the caller is filled & \code{None} & \code{str}; \code{aux} is \code{bytes}\\
C++ & \code{wfdbx::\allowbreak isigopen(...)} (namespace \code{wfdbx}, header
\file{<wfdbx/wfdb.hpp>}) & as in C: arrays and structures allocated by the
caller, passed by pointer & \code{nullptr} & \code{std::string}, or
\code{std::optional<\allowbreak std::string>} where C may return
\code{NULL}\\
Rust & \code{wfdbx::\allowbreak isigopen(...)}, or the method \code{w.isigopen(...)} of a
\code{Wfdb} value & a mutable slice or reference passed by the caller is
filled & \code{None} (\code{wfdbx::NULL}) & \code{Vec<u8>}, or
\code{Option<Vec<u8>>} where C may return \code{NULL}\\
Julia & \code{isigopen(...)} after \code{using WFDBX}, or
\code{WFDBX.\allowbreak isigopen(...)}; optionally a \code{Wfdb} object first & a vector
or mutable structure passed by the caller is filled; a \code{Ref} receives a
pointer result & \code{nothing} & \code{Vector\{UInt8\}}; \code{String} for
\code{timstr}, \code{mstimstr}, \code{datstr}\\
R & \code{isigopen(...)} (package \code{wfdbx}) & returned with the C result
in a list (\code{r\$nsig}, \code{r\$si}; \code{g\$stat}, \code{g\$v}) & \code{NULL}
& character strings; \code{aux} is a raw vector\\
MATLAB/ Octave & \code{wfdbx.isigopen(...)} (package \code{+wfdbx}) & returned as
extra output values (\code{[nsig, si] = ...}) & \code{[]} & \code{char} row
vectors; \code{aux} is \code{uint8}\\
Kotlin & \code{isigopen(...)} (package \code{wfdbx}), or the extension
function \code{w.isigopen(...)} of a \code{Wfdb} object & an array or object
passed by the caller is filled & \code{null} & \code{ByteArray}, or
\code{ByteArray?} where C may return \code{NULL}\\
\end{xltabular}

\subsection{Python}
\label{sec:library-python}

The functions, the structures and the constants are in the package
\code{wfdbx} (\code{import wfdbx}). The C structures are the classes
\code{Siginfo}, \code{Anninfo}, \code{Annotation}, \code{Calinfo} and
\code{Seginfo}, with the C field names; their constructors take the fields
as keyword arguments, with C's zero values, except that \code{Siginfo.spf} is
1 by default.

\begin{itemize}
\item Where C fills an array of structures (\code{isigopen},
  \code{osigopen}), the caller passes a list, which is emptied and receives
  one \code{Siginfo} per signal. Where C fills an array of samples
  (\code{getvec}, \code{getframe}), the caller passes a list with room for
  every value (\code{[0] * nsig}), whose first elements are replaced.
  \code{putvec} reads a list (or any sequence).
\item Where C fills a structure (\code{getann}, \code{getcal}), the caller
  passes an object (\code{Annotation()}, \code{Calinfo()}), whose fields are
  replaced.
\item \code{getseginfo()} takes no argument and returns the list of
  \code{Seginfo} objects (empty if the record is not a multi-segment
  record), instead of a count and a pointer.
\item The counts that only give the length of an array have defaults:
  \code{annopen(record, aiarray)}, \code{osigfopen(siarray)},
  \code{setheader(record, siarray)} and \code{setmsheader(record, names)}
  use the length of the list; \code{isigopen} and \code{osigopen} take
  \code{siarray=None, nsig=0}, \code{getinfo} and \code{calopen} take
  \code{None}.
\item Strings are \code{str}; an annotation's \code{aux} is \code{None} or
  \code{bytes}, and \code{Annotation.aux\_string()} returns its text.
  \code{wfdb\_error(fmt, *args)} formats as C's \code{printf}.
\item \code{getframes(n)} reads up to \code{n} sample vectors into a NumPy
  array (\cref{sec:library-additions}).
\item The functions of \file{<wfdb/ecgmap.h>} (\code{isann}, \code{isqrs},
  \code{map1}, ...) and the annotation codes (\code{NORMAL}, \code{PVC}, ...)
  are in the package; \code{isann} returns \code{True} or \code{False}
  (C: 1 or 0).
\end{itemize}

\begin{lstlisting}[language=Python,style=example]
import wfdbx

si = []
nsig = wfdbx.isigopen("100s", si, 2)    # si receives 2 Siginfo objects
v = [0] * 2
if nsig == 2 and wfdbx.getvec(v) > 0:   # v receives one sample per signal
    print(si[0].desc, v[0], v[1])       # MLII 995 1011
wfdbx.wfdbquit()
\end{lstlisting}

\subsection{C++}
\label{sec:library-cpp}

\file{<wfdbx/wfdb.hpp>} declares the functions, types, structures and
constants in namespace \code{wfdbx}; \file{<wfdbx/ecgcodes.hpp>} defines the
annotation codes, and \file{<wfdbx/ecgmap.hpp>} the macros of
\file{<wfdb/ecgmap.h>} as inline functions (\code{isann}, \code{isqrs}, ...,
which call \code{wfdb\_isann}, \code{wfdb\_isqrs}, ...).

\begin{itemize}
\item Arrays and structures are allocated by the caller and passed by
  pointer, as in C. String arguments are \code{const char *}, and
  \code{nullptr} is accepted where C accepts \code{NULL}.
\item Strings are returned by value, as \code{std::string}, or as
  \code{std::optional<std::string>} where the C function may return
  \code{NULL} (\code{annstr}, \code{ecgstr}, \code{anndesc}, \code{getinfo},
  \code{wfdbfile}); unlike C's pointers to static buffers, they stay valid
  after later calls.
\item The string fields of \code{WFDB\_Siginfo}, \code{WFDB\_Calinfo} and
  \code{WFDB\_Annotation} (including \code{aux}) are
  \code{std::optional<std::string>}, since C distinguishes \code{NULL} from
  \code{""}; \code{WFDB\_Anninfo::name} is a \code{std::string}. The
  structures are also available under the shorter names \code{Siginfo},
  \code{Anninfo}, \code{Annotation}, \code{Calinfo} and \code{Seginfo}.
\item \code{WFDB\_Time} is \code{long long} (64 bits); \code{getseginfo}
  takes a \code{const WFDB\_Seginfo **}.
\item \code{getvecs} reads many sample vectors at once
  (\cref{sec:library-additions}). The library is not thread-safe, as in C.
\end{itemize}

\begin{lstlisting}[language=C++,style=example]
#include <cstdio>
#include <wfdbx/wfdb.hpp>
using namespace wfdbx;

int main() {
    WFDB_Siginfo si[2];
    WFDB_Sample v[2];
    int nsig = isigopen("100s", si, 2);
    if (nsig == 2 && getvec(v) > 0)
        std::printf("%s %d %d\n", si[0].desc->c_str(), v[0], v[1]);
    wfdbquit();
    return 0;
}
\end{lstlisting}

\subsection{Rust}
\label{sec:library-rust}

The state of the library is held by a \code{Wfdb} value, and each function is
a method of it (\code{w.isigopen(...)}). Every function is also a free
function of the crate with the same name and arguments
(\code{wfdbx::\allowbreak isigopen(...)}), which uses a default \code{Wfdb} value, one per
thread; separate \code{Wfdb} values share no state, so that several records
can be open at once. The functions that use no state of their own are free
functions only: \code{wfdbldflags}, \code{wfdbcflags}, \code{wfdbdefwfdb},
\code{wfdbdefwfdbcal}, the count of write errors (\code{wfdbx\_write\_errors},
\code{wfdbx\_note\_write\_error}, which belongs to the process),
\code{isann}, \code{ammap} and \code{mamap} (and their \code{wfdb\_} forms),
and the macros of \file{<wfdb/ecgmap.h>} that use the mapping tables
(\code{isqrs}, \code{map1}, ...), which use those of the default value (the
methods \code{w.wfdb\_isqrs(...)}, \code{w.wfdb\_map1(...)}, ... use those of
\code{w}).

\begin{itemize}
\item Where C fills an array, the caller passes a mutable slice
  (\code{\&mut si}, \code{\&mut v}); where C fills a structure, a mutable
  reference (\code{\&mut annot}). C's \code{NULL} array, as in
  \code{isigopen(record, NULL, 0)}, is an empty slice.
\item String arguments accept \code{\&str}, \code{\&[u8]} or \code{b"..."};
  where C accepts \code{NULL}, they also accept \code{Option}s, and
  \code{wfdbx::NULL} is a typed \code{None}. Strings are returned as
  \code{Vec<u8>}, or \code{Option<Vec<u8>>} where C may return \code{NULL}
  (\code{anndesc}, \code{getinfo}, \code{wfdbfile}, ...).
\item C's \code{int} is \code{i32}, \code{long} and \code{WFDB\_Time} are
  \code{i64}, \code{WFDB\_Signal}, \code{WFDB\_Group} and
  \code{WFDB\_Annotator} are \code{u32}. The structures (\code{Siginfo},
  \code{Anninfo}, \code{Annotation}, \code{Calinfo}, \code{Seginfo}) have the
  C field names; \code{Default::default()} gives C's zero values; the
  \code{char} fields of \code{Annotation} are \code{i8} (\code{chan}:
  \code{u8}).
\item \code{getseginfo(\&mut segs)} takes an \code{\&mut Option<Vec<Seginfo>>},
  which receives a copy of the segment table (\code{None} if there is none).
\item \code{wfdb\_error(msg)} takes a message already formatted (with the
  crate's \code{cformat!} macro, which formats as C's \code{printf}), and
  \code{wfdb\_error\_fmt(fmt, args)} formats it.
\item \code{getvecs} and \code{getvecs\_status} read many sample vectors at
  once (\cref{sec:library-additions}); \code{wfdbx\_fclose\_output} closes a
  file that a program wrote itself and counts a write error
  (\cref{sec:library-additions}).
\end{itemize}

\begin{lstlisting}[language=Rust,style=example]
use wfdbx::Siginfo;

fn main() {
    let mut si = vec![Siginfo::default(); 2];
    let nsig = wfdbx::isigopen("100s", &mut si, 2);
    let mut v = vec![0i32; 2];
    if nsig == 2 && wfdbx::getvec(&mut v) > 0 {
        let desc = si[0].desc.clone().unwrap_or_default();
        println!("{} {} {}", String::from_utf8_lossy(&desc), v[0], v[1]);
    }
    wfdbx::wfdbquit();
}
\end{lstlisting}

\subsection{Julia}
\label{sec:library-julia}

The functions are in the module \code{WFDBX}. \code{using WFDBX} makes the
names of most of them available; those that it does not export are called
with the module prefix: \code{WFDBX.wfdb\_error}, \code{WFDBX.wfdb\_me\_fatal},
and the function forms of the macros of \file{<wfdb/ecgmap.h>},
\code{WFDBX.wfdb\_isann}, \code{WFDBX.wfdb\_isqrs}, ... (\code{isann},
\code{isqrs}, ... are exported). The state of the library is held by a
\code{Wfdb} object: every function that uses it takes one as an optional
first argument (\code{isigopen(w, "100s", si, 2)}) and otherwise uses the
default object, \code{lib()}; \code{Wfdb()} creates an independent one. The
functions that use no such state take none: \code{wfdbldflags},
\code{wfdbcflags}, \code{wfdbdefwfdb}, \code{wfdbdefwfdbcal}, \code{isann},
\code{ammap}, \code{mamap} (and their \code{wfdb\_} forms), and
\code{wfdbx\_write\_errors} and \code{wfdbx\_note\_write\_error}, whose
count belongs to the process.

\begin{itemize}
\item Where C fills an array, the caller passes a vector: of \code{Siginfo}
  structures (\code{[Siginfo() for \_ in 1:n]}), or of samples, which must
  be a \code{Vector\{Int32\}}. Where C fills a structure, the caller passes
  a mutable structure (\code{Annotation()}). \code{getseginfo(r)} takes a
  \code{Ref}, which receives a vector of \code{Seginfo} (or \code{nothing}).
\item C's \code{NULL} is \code{nothing}. Arguments may be \code{String}s or
  bytes. Text is returned as bytes (\code{Vector\{UInt8\}}; \code{String(x)}
  converts it), except the results of \code{timstr}, \code{mstimstr} and
  \code{datstr}, which are \code{String}s; the string fields of the
  structures are bytes, or \code{nothing}.
\item The integer types are those of C: \code{WFDB\_Sample} and \code{int} are
  \code{Int32}, \code{WFDB\_Time} and \code{long} are \code{Int64}; the
  \code{char} fields of \code{Annotation} are \code{Int8} (\code{chan}:
  \code{UInt8}).
\item \code{wfdb\_error(fmt, args...)} formats as C's \code{printf}.
  \code{wfdbx\_legacy\_args(args)} takes the arguments without the program
  name (\cref{sec:library-additions}).
\item \code{getvecs} and \code{getvecs\_status} read many sample vectors at
  once (\cref{sec:library-additions}).
\end{itemize}

\begin{lstlisting}[language=Julia,style=example]
using WFDBX

si = [Siginfo() for _ in 1:2]
nsig = isigopen("100s", si, 2)          # si receives the specifications
v = zeros(Int32, 2)
if nsig == 2 && getvec(v) > 0           # v receives one sample per signal
    println(String(si[1].desc), " ", v[1], " ", v[2])
end
wfdbquit()
\end{lstlisting}

\subsection{R}
\label{sec:library-r}

The functions are those of the package \code{wfdbx} (\code{library(wfdbx)}).
R has no output arguments, so a function that fills an array or a
structure in C returns it, together with the C return value, in a list; the
argument that C fills is omitted.

\begin{itemize}
\item \code{isigopen(record, nsig)} and \code{osigopen(record, nsig)} return
  \code{list(nsig, si)}: the C result and the list of \code{siginfo} lists;
  \code{wfdbinit(record, aiarray, nsig)} (without \code{nann}) returns
  \code{list(stat, si)}; \code{getvec()} and \code{getframe()} return
  \code{list(stat, v)}; \code{getann(an)} returns \code{list(stat, annot)};
  \code{getcal(desc, units)} returns \code{list(stat, cal)};
  \code{setheader} returns \code{list(stat, si)} (with the units and
  descriptions as \code{setheader} truncates them);
  \code{getseginfo()} returns the list of \code{seginfo} lists;
  \code{wfdbx\_legacy\_args(argv)} returns \code{list(argv, legacy)}.
\item The C structures are lists with the C field names, made by
  \code{siginfo()}, \code{anninfo()}, \code{annotation()}, \code{calinfo()}
  and \code{seginfo()} (\code{siginfo()} gives \code{spf} 1). Numbers are
  doubles; strings are character strings, and C's \code{NULL} is
  \code{NULL}; an annotation's \code{aux} is \code{NULL} or a raw vector,
  and \code{aux\_string(aux)} returns its text.
\item The C function \code{sample} is \code{wfdb\_sample}, so that it does not
  mask R's \code{sample}. The functions of \file{<wfdb/ecgmap.h>} accept
  vectors of codes. \code{wfdb\_error(fmt, ...)} formats as C's
  \code{printf}.
\item \code{getframes(n)} reads up to \code{n} sample vectors into a matrix
  (\cref{sec:library-additions}); \code{putframes(m)} writes the rows of a
  matrix with \code{putvec}.
\end{itemize}

\begin{lstlisting}[language=R,style=example]
library(wfdbx)

r <- isigopen("100s", 2)           # r$nsig, and r$si: the specifications
g <- getvec()                      # g$stat, and g$v: one sample per signal
if (r$nsig == 2 && g$stat > 0)
  cat(r$si[[1]]$desc, g$v, "\n")
wfdbquit()
\end{lstlisting}

\subsection{MATLAB/Octave}
\label{sec:library-matlab}

The functions are those of the package \file{+wfdbx}, called with the prefix
\code{wfdbx.} (\code{wfdbx.isigopen(...)}). An array or structure that C
fills is returned as an additional output value, and omitted from the
arguments.

\begin{itemize}
\item \code{[nsig, si] = wfdbx.isigopen(record, nsig)} and
  \code{[nsig, si] = wfdbx.osigopen(record, nsig)} return the C result and a
  struct array of \code{wfdbx.siginfo} structures; \code{[stat, si] =
  wfdbx.wfdbinit(record, aiarray, nann, nsig)}; \code{[stat, v] =
  wfdbx.getvec()} and \code{[stat, v] = wfdbx.getframe()} (an initial vector
  may be passed: its elements are kept where C leaves them unchanged);
  \code{[stat, annot] = wfdbx.getann(an)}; \code{[stat, cal] =
  wfdbx.getcal(desc, units)}; \code{[stat, si] = wfdbx.setheader(record, si,
  nsig)}; \code{segs = wfdbx.getseginfo()} returns a struct array;
  \code{[r, argv] = wfdbx.wfdbx\_legacy\_args(argv)}.
\item The C structures are MATLAB structures with the C field names, made by
  \code{wfdbx.siginfo(n)}, \code{wfdbx.anninfo(name, stat)},
  \code{wfdbx.annotation()}, \code{wfdbx.calinfo()} and \code{wfdbx.seginfo(n)}
  (\code{siginfo} gives \code{spf} 1). Samples, times and counts are
  doubles. Strings are \code{char} row vectors (record names can also be
  MATLAB strings), and C's \code{NULL} is \code{[]}; an annotation's
  \code{aux} is \code{[]} or a \code{uint8} row vector.
\item Constants are functions without arguments: \code{wfdbx.WFDB\_READ()},
  \code{wfdbx.NORMAL()}, ...; \code{wfdbx.ecgcodes()} returns all the
  annotation codes in one structure. The calibration constants
  (\code{WFDB\_AC\_COUPLED}, ...) and \code{WFDB\_FMT\_LIST} are not defined;
  their values are given in \cref{tab:library-constants}.
\item \code{wfdbx.wfdb\_error(fmt, ...)} formats as C's \code{printf}.
  \code{wfdbx.getframes(n)} reads up to \code{n} sample vectors into a matrix
  (\cref{sec:library-additions}), and \code{wfdbx.putframes(M)} writes the
  rows of a matrix.
\end{itemize}

\begin{lstlisting}[language=Matlab,style=example]
[nsig, si] = wfdbx.isigopen('100s', 2);   % si: the specifications
[stat, v] = wfdbx.getvec();               % v: one sample per signal
if nsig == 2 && stat > 0
    fprintf('%s %d %d\n', si(1).desc, v(1), v(2));
end
wfdbx.wfdbquit();
\end{lstlisting}

\subsection{Kotlin}
\label{sec:library-kotlin}

The functions are in the package \code{wfdbx} (\code{import wfdbx.*}). The
state of the library is held by a \code{Wfdb} object, and each function is an
extension function of it (\code{w.isigopen(...)}); every function is also a
top-level function with the same name and arguments
(\code{isigopen(...)}), which uses a default \code{Wfdb} object, one per
thread. Separate \code{Wfdb} objects share no state; one object must be used
by one thread at a time. As in Rust, \code{wfdbldflags}, \code{wfdbcflags},
\code{wfdbdefwfdb}, \code{wfdbdefwfdbcal}, \code{wfdbx\_write\_errors},
\code{wfdbx\_note\_write\_error}, \code{isann}, \code{ammap},
\code{mamap} (and their \code{wfdb\_} forms) and the macros of
\file{<wfdb/ecgmap.h>} that use the mapping tables (\code{isqrs},
\code{map1}, ..., on the default object; the extension functions are
\code{w.wfdb\_isqrs(...)}, ...) are top-level functions only.

\begin{itemize}
\item Where C fills an array, the caller passes an array
  (\code{Array<Siginfo>}, \code{IntArray}); where C fills a structure, an
  object (\code{Annotation()}). \code{getseginfo(list)} takes a
  \code{MutableList<Seginfo>?}, which receives the segments.
\item Strings are bytes (\code{ByteArray}), as in C; most functions also
  accept a \code{String} (one character per byte), and \code{latin1()} (in
  \code{wfdbx.cstd}) converts in either direction. C's \code{NULL} is
  \code{null}; results that C may return as \code{NULL} are
  \code{ByteArray?}.
\item C's \code{int} and \code{WFDB\_Sample} are \code{Int}, \code{long} and
  \code{WFDB\_Time} are \code{Long}, \code{double} is \code{Double}. The
  structures (\code{Siginfo}, \code{Anninfo}, \code{Annotation},
  \code{Calinfo}, \code{Seginfo}) have the C field names and C's zero values;
  the \code{char} fields of \code{Annotation} are \code{Int} properties that
  keep C's range (\code{anntyp}, \code{subtyp}, \code{num}: -128 to 127;
  \code{chan}: 0 to 255).
\item \code{wfdb\_error(msg)} takes a message already formatted (with
  \code{cformat} of \code{wfdbx.cstd}, which formats as C's \code{printf}).
\item \code{getvecs} and \code{getvecs\_status} read many sample vectors at
  once, and \code{wfdbx\_fclose\_output} closes a file that a program wrote
  itself (\cref{sec:library-additions}).
\end{itemize}

\begin{lstlisting}[language=Kotlin,style=example]
import wfdbx.*
import wfdbx.cstd.latin1

fun main() {
    val si = Array(2) { Siginfo() }
    val nsig = isigopen("100s", si, 2)
    val v = IntArray(2)
    if (nsig == 2 && getvec(v) > 0)
        println("${si[0].desc!!.latin1()} ${v[0]} ${v[1]}")
    wfdbquit()
}
\end{lstlisting}

\begin{note}[Functions that exist only in some languages]
Every function documented in this chapter exists in the seven languages
(\file{docs/INVENTORY.md}). The reader of many sample vectors has two names:
\code{getframes} in Python, R and MATLAB/Octave, and \code{getvecs} (with
\code{getvecs\_status}) in C++, Rust, Julia and Kotlin; it is documented under
\code{getframes}. A few helpers have no counterpart in the Python
implementation, and therefore no entry here: \code{putframes} (R,
MATLAB/Octave), \code{wfdbx\_fclose\_output} (Rust, Kotlin),
\code{wfdb\_error\_fmt} (Rust), \code{aux\_string} (R; a method of
\code{Annotation} in Python) and the structure constructors of R and
MATLAB/Octave. In R, \code{sample} is named \code{wfdb\_sample}.
\end{note}

% ==========================================================================
\section{Opening and closing records}
\label{sec:library-open}

A record is opened for reading by \code{isigopen} (its signals) and
\code{annopen} (its annotation files), or by \code{wfdbinit}, which calls
both; output signals are opened by \code{osigopen} or \code{osigfopen}, and
output annotation files by \code{annopen}. Programs that write files must call
\code{wfdbquit} (or close the files with the functions given in each entry)
before they exit.

A record name is found as \cref{chap:behavior} describes: a relative name
under the input directory (for reading) or the output directory (for
writing), an absolute name, a \code{\textasciitilde} name or a URL as given; a
name that begins with \code{+} keeps the files of the same kind already open,
and the record name \code{-} is the standard input or output. A trailing
\code{.hea} is removed from a record name.

\apientry{annopen}{Open annotation files}
\begin{lstlisting}[language=C,style=signature]
int annopen(char *record, const WFDB_Anninfo *aiarray, unsigned int nann);
\end{lstlisting}
Opens, for \code{record}, the \code{nann} annotation files described by
\code{aiarray}: for reading if their \code{stat} is \code{WFDB\_READ} or
\code{WFDB\_AHA\_READ}, for writing (created, or truncated) if it is
\code{WFDB\_WRITE} or \code{WFDB\_AHA\_WRITE}. The file of annotator
\code{atr} of record \code{100s} is \file{100s.atr}. Input and output
annotators are numbered separately, from 0, in the order of \code{aiarray}:
these are the numbers given to \code{getann} and \code{putann}. All the
annotation files already open are closed first (output files are completed,
and sorted if needed, as by \code{oannclose}), unless \code{record} begins with
\code{+}, which is then removed. All the \code{stat} values are checked before
any file is opened. An input file whose contents show the other format (MIT or
AHA) than its \code{stat} says is read in the format it has, with a warning.
The modification labels (mnemonics and descriptions defined at the beginning
of an input file) are applied to the tables of \code{annstr} and
\code{anndesc}; the time resolution recorded in an input file is used by
\code{getann} (see \code{getiaorigfreq}). An output annotator name may contain
only letters, digits, \code{\_}, \code{\textasciitilde}, \code{-} and
\code{/}.

With \code{nann} 0, \code{annopen} only closes the open annotation files
(\code{annopen("", NULL, 0)}) and reports whether that succeeded.
\begin{params}
\item[record] The record name.
\item[aiarray] An array of \code{nann} \code{WFDB\_Anninfo} structures
  (\cref{sec:library-structures}): the annotator name (\code{name}) and the
  mode (\code{stat}) of each file.
\item[nann] The number of annotators in \code{aiarray}.
\end{params}
\returns{0 on success; -3 if an input file cannot be opened; -4 if an output
file cannot be created or its annotator name is illegal; -5 if a \code{stat}
is not one of the four values. With \code{nann} 0: 0, -6 if an output file
closed by the call was out of order (whether or not it could be sorted), or
-7 if one could not be written completely.}
\raises{\code{annopen: can't read annotator ATR for record REC};
\code{annopen: can't write annotator ATR for record REC};
\code{annopen: illegal stat N for annotator ATR, record REC}; \code{init:
illegal character N in annotator name}.}

WFDB-x: input files are read only where their name says, output files are
created under the output directory, and out-of-order output files closed by
the call are sorted in the same process (\cref{chap:behavior}).

\apientry{isigopen}{Open the input signals of a record}
\begin{lstlisting}[language=C,style=signature]
int isigopen(char *record, WFDB_Siginfo *siarray, int nsig);
\end{lstlisting}
Reads the header of \code{record} and opens its signal files for reading.
The input signals already open are closed first, unless \code{record} begins
with \code{+}, which is then removed (the new signals are numbered after
those open). The header also sets the sampling and counter frequencies, the
base counter value, the base time and date, and the length of the record,
which the time conversion functions use.

With \code{nsig} greater than 0, at most \code{nsig} signals are opened, and
the structures \code{siarray[0]} to \code{siarray[n-1]} are filled with their
specifications (signal numbers are assigned in the order of the header).
Signals stored in the same file (a signal group) are opened together: a
group is skipped if its file cannot be opened, or if its signals do not fit
in the remaining entries (a record whose signals are all in one file is
therefore not opened at all if \code{nsig} is smaller than its number of
signals).

With \code{nsig} 0 or negative, no signals are opened (and the open input
signals are closed, as above): the header is read, up to \code{-nsig}
structures of \code{siarray} are filled from it (\code{siarray} may be
\code{NULL} when \code{nsig} is 0), and the number of signals named in the
header is returned. \code{isigopen(record, NULL, 0)} thus gives the number of
signals, to allocate \code{siarray} before opening them, and sets the time
base of the record for \code{strtim} and \code{timstr}.

A record name that contains a \code{.}, other than in a final \code{.hea},
names an EDF file, read as a record (\cref{chap:formats}). A multi-segment
record opens its segments as they are read (see \code{getseginfo}).
\begin{params}
\item[record] The record name.
\item[siarray] An array of at least \code{nsig} (or \code{-nsig})
  \code{WFDB\_Siginfo} structures, filled by the function.
\item[nsig] The maximum number of signals to open; 0 or a negative number to
  read the header only.
\end{params}
\returns{With \code{nsig} > 0: the number of signals opened (possibly fewer
than \code{nsig}), 0 if the record has no signals or none could be opened.
With \code{nsig} <= 0: the number of signals of the record. In either case:
-1 if the header cannot be read (probably an incorrect record name), -2 if it
is not a valid header, -3 if the internal limits are exceeded (the total
number of samples per frame, or the skew, too large), in which case all the
input signals are closed.}
\raises{\code{init: can't open header for record REC}; \code{isigopen:
record REC has no signals}; \code{isigopen: none of the signals for record
REC is readable}; the messages of the header parser (\code{init: ...}).}

WFDB-x (\cref{chap:behavior}): the header is read only where the record name
says, and the signal files and segments that it names only from the header's
directory (\code{init: file F named in the header of record REC is not in
the header's directory} for a name that leaves it); a header without a
sampling frequency, ADC resolution or units leaves them unspecified (0, 0,
\code{NULL}) instead of assuming 250 Hz, a resolution and millivolts; a header
with a negative or impossibly large number of signals gives \code{init:
unexpected EOF in header file for record REC} (C tries to allocate memory for
that count).

\apientry{osigopen}{Open output signals as described by a header}
\begin{lstlisting}[language=C,style=signature]
int osigopen(char *record, WFDB_Siginfo *siarray, unsigned int nsig);
\end{lstlisting}
Reads the header of the existing record \code{record} and opens for writing
(creating or truncating them) the signal files of its first \code{nsig}
signals, with the specifications of the header; \code{siarray} receives them
(with \code{nsamp} and \code{cksum} 0). The output signals already open are
closed first, unless \code{record} begins with \code{+}, which is then
removed. The samples are then written with \code{putvec}, and a header for
the new signals with \code{newheader}. A header of an existing database
record names that record's signal files, which would be overwritten: the
header used here is usually written for the purpose (a ``piped'' record such
as \code{8l}, whose signal files are \file{data0} and \file{data1}).
\begin{params}
\item[record] The record whose header specifies the output signals.
\item[siarray] An array of at least \code{nsig} \code{WFDB\_Siginfo}
  structures, filled by the function.
\item[nsig] The number of output signals to open.
\end{params}
\returns{The number of output signals opened (\code{nsig}, 0 when \code{nsig}
is 0); -1 if the header cannot be read; -2 if it is not valid, or if the
signals of a group differ in format or block size; -3 if the record has fewer
than \code{nsig} signals or a signal file cannot be created.}
\raises{\code{osigopen: record REC has fewer signals than needed};
\code{osigopen: can't open FILE}; \code{osigopen: error in specification of
signal N or M}.}

The Programmer's Guide allows \code{nsig} 0 with \code{siarray} \code{NULL},
to close the output signals; the C function then reads before the start of
its internal array. In WFDB-x, in every language, \code{osigopen(record, NULL,
0)} closes the output signals already open (unless \code{record} begins with
\code{+}), reads the header (and returns -1 if it cannot be read), opens
nothing, and returns 0 (\cref{sec:crashes}). \code{osigfopen(NULL, 0)} also
closes the output signals.

WFDB-x: a signal file is created under the output directory, and a signal
file that is one of the signal files of the open input record is refused
(\code{osigopen: output file FILE is also an input file}, -3; legacy mode does
not check); see \cref{chap:behavior}.

\apientry{osigfopen}{Open output signals as described by the caller}
\begin{lstlisting}[language=C,style=signature]
int osigfopen(const WFDB_Siginfo *siarray, unsigned int nsig);
\end{lstlisting}
Closes the open output signals, and opens \code{nsig} output signals as
\code{siarray} describes them: file name, description, units, gain, format,
samples per frame, block size, ADC resolution and zero, baseline
(\code{initval}, \code{nsamp} and \code{cksum} are not used). Signals written
to the same file (a multiplexed signal file) must be consecutive and have the
same \code{fname} and \code{group}, and the same format and block size; each
new file has the next group number, starting at 0. A \code{fname} of
\code{-} is the standard output; a signal of format 0 with no file name (or
\code{\textasciitilde}) is not written, and a signal of format 0 with a file
name is written in format 16. An \code{spf} smaller than 1 is taken as 1.

\code{osigfopen(NULL, 0)} closes the output signals (and the header opened by
\code{newheader}, \code{setheader} or \code{setmsheader}) and reports whether
they were written successfully.
\begin{params}
\item[siarray] An array of \code{nsig} \code{WFDB\_Siginfo} structures,
  filled by the caller.
\item[nsig] The number of output signals.
\end{params}
\returns{The number of output signals opened (\code{nsig}); -2 for an error in
a specification (\code{fname} and \code{desc} together longer than 200
characters, a negative \code{bsize}, a format not in
\code{WFDB\_FMT\_LIST}, a wrong group or file name, or signals of a group with
different formats or block sizes); -3 if a signal file cannot be created.
With \code{nsig} 0: 0, or -4 if a signal file or header closed by the call
could not be written completely.}
\raises{\code{osigfopen: error in specification of signal N};
\code{osigfopen: incorrect file name or group for signal N};
\code{osigfopen: can't open FILE}; \code{osigclose: write error in signal
group N}.}

WFDB-x: as for \code{osigopen}, files are created under the output directory,
and a signal file of the open input record is refused (\code{osigfopen:
output file FILE is also an input file}); see \cref{chap:behavior}.

\apientry{wfdbinit}{Open the annotation files and the input signals}
\begin{lstlisting}[language=C,style=signature]
int wfdbinit(char *record, const WFDB_Anninfo *aiarray, unsigned int nann,
             WFDB_Siginfo *siarray, unsigned int nsig);
\end{lstlisting}
Calls \code{annopen(record, aiarray, nann)} and, if it returns 0,
\code{isigopen(record, siarray, nsig)}. The Programmer's Guide advises
against using \code{wfdbinit} together with \code{setifreq}: \code{setifreq}
must be called after the signals are opened and before the annotation files
are, to apply to the annotation times.
\begin{params}
\item[record] The record name.
\item[aiarray] The annotators, as for \code{annopen}.
\item[nann] The number of annotators.
\item[siarray] An array of at least \code{nsig} \code{WFDB\_Siginfo}
  structures, as for \code{isigopen}.
\item[nsig] The maximum number of input signals to open.
\end{params}
\returns{The negative result of \code{annopen} if it fails; otherwise the
result of \code{isigopen}: the number of signals opened, 0 if none (no
message if \code{nsig} is 0), or -1, -2, -3.}

In R, \code{wfdbinit(record, aiarray, nsig)} has no \code{nann} argument, and
returns \code{list(stat, si)}; in MATLAB/Octave, \code{[stat, si] =
wfdbx.wfdbinit(record, aiarray, nann, nsig)}.

\apientry{wfdbquit}{Close all files and reset the library}
\begin{lstlisting}[language=C,style=signature]
void wfdbquit(void);
\end{lstlisting}
Closes all the annotation files (completing and, if needed, sorting the
output files, as \code{oannclose}), the info file opened by \code{setinfo},
the input and output signal files (writing out the buffered samples) and the
header being written, and releases the buffers of \code{sample} and the info
strings of \code{getinfo}. It resets the sampling frequencies, the base time
and date, the record length and the segments of the open record. It does not
write a header (\code{newheader} or \code{setheader} must be called first),
and, as in C, it does not reset the counter frequency and base counter value,
the output annotation time resolution (\code{setafreq}), the buffer sizes
(\code{setibsize}, \code{setobsize}), the calibration list, or the tables of
mnemonics, descriptions and code mappings (\cref{chap:behavior}). Programs
that write signals or annotations must call it (or the closing functions of
each kind of file) before they exit.

WFDB-x: the input directory, the output directory and the network setting
are kept; in legacy mode, the WFDB path is restored as by \code{resetwfdb}
(\cref{chap:behavior}).

\apientry{wfdbflush}{Write the buffered output}
\begin{lstlisting}[language=C,style=signature]
void wfdbflush(void);
\end{lstlisting}
Writes the buffered output annotations and samples to their files, so that
the files are up to date while they stay open.

% ==========================================================================
\section{Reading and writing signals}
\label{sec:library-signals}

The input signals are read one sample vector at a time, sequentially, with
\code{getvec} (one sample of each signal) or \code{getframe} (every sample of
a frame); \code{isigsettime}, \code{isgsettime} and \code{tnextvec} move to
another time, and \code{sample} gives random access. Samples are integers in
ADC units; a missing or invalid sample is \code{WFDB\_INVALID\_SAMPLE}
(-32768) unless the \code{WFDB\_GVPAD} mode is set. Times are in sample
intervals: of the frame, of the highest sampling frequency in
\code{WFDB\_HIGHRES} mode, or of the frequency set by \code{setifreq}.

\apientry{getvec}{Read a sample of each input signal}
\begin{lstlisting}[language=C,style=signature]
int getvec(WFDB_Sample *vector);
\end{lstlisting}
Reads the next sample of each input signal into \code{vector}: \code{vector[i]}
receives the sample of signal \code{i}. In \code{WFDB\_LOWRES} mode (the
default), a signal with several samples per frame gives the mean of the
samples of the frame (rounded to the nearest integer, halves upward; invalid
if one of them is); in \code{WFDB\_HIGHRES} mode, each call gives the next
sample at the highest sampling frequency, and the other signals are repeated
(see \code{setgvmode}). After \code{setifreq}, the signals are resampled at
the chosen frequency by linear interpolation. Skewed signals are corrected
for their skew. The checksum of a signal is verified when the number of
samples given in the header has been read, unless the signals were
positioned (\code{isigsettime}, \code{isgsettime}) before.
\begin{params}
\item[vector] An array with room for one sample per input signal, filled by
  the function.
\end{params}
\returns{The number of input signals, on success; 0 if no input signals are
open; -1 at the end of the data (the contents of \code{vector} are then not
valid); -3 at an unexpected end of a signal file (or a FLAC decoding error);
-4 on a checksum error, detected on the last frame (which is valid).}
\raises{\code{getvec: unexpected EOF in signal N}; \code{getvec: checksum
error in signal N}; \code{getvec: error opening segment SEG}.}

\apientry{getframe}{Read a frame of the input signals}
\begin{lstlisting}[language=C,style=signature]
int getframe(WFDB_Sample *vector);
\end{lstlisting}
Reads the next frame: every sample of every input signal for one frame
interval, the \code{spf} samples of signal 0 first, then those of signal 1,
and so on. For a record whose signals all have one sample per frame,
\code{getframe} and \code{getvec} (without \code{setifreq}) read the same
vectors. \code{setgvmode} and \code{setifreq} do not affect it.
\begin{params}
\item[vector] An array with room for the sum of the \code{spf} fields of the
  input signals, filled by the function.
\end{params}
\returns{As \code{getvec}.}

\apientry{putvec}{Write a sample to each output signal}
\begin{lstlisting}[language=C,style=signature]
int putvec(const WFDB_Sample *vector);
\end{lstlisting}
Writes a frame to the output signals: \code{vector} holds the \code{spf}
samples of output signal 0, then those of signal 1, and so on (one sample per
signal when every \code{spf} is 1). The number of samples and the checksum of
each signal are kept, for \code{newheader}. A sample equal to
\code{WFDB\_INVALID\_SAMPLE} is written as the invalid-sample value of the
format; a sample outside the range of the format is written as C writes it
(its low-order bits). In format 8 (first differences), a difference outside
-128 to 127 is limited to that range, and the following samples correct the
level as soon as the differences allow.
\begin{params}
\item[vector] The samples, given by the caller.
\end{params}
\returns{The number of output signals on success; 0 if a difference was too
large in format 8; -1 on a write error.}
\raises{\code{putvec: write error in signal N}.}

\apientry{isigsettime}{Skip to a time in all input signals}
\begin{lstlisting}[language=C,style=signature]
int isigsettime(WFDB_Time t);
\end{lstlisting}
Positions the input signals so that the next sample read is the one at time
\code{t}; only the magnitude of \code{t} counts, so that a value returned by
\code{strtim} for a time of day can be given as it is. Nothing is done if the
signals are already at \code{t} (without \code{setifreq}). A pipe, such as the
standard input, cannot be positioned backward.
\begin{params}
\item[t] The time, in sample intervals (see the introduction of this
  section).
\end{params}
\returns{0 on success (also if no input signals are open); -1 at the end of
the data or on an improper seek; -2 if a signal group is incorrect.}
\raises{\code{isigsettime: improper seek on signal group N}.}

\apientry{isgsettime}{Skip to a time in one signal group}
\begin{lstlisting}[language=C,style=signature]
int isgsettime(WFDB_Group sgroup, WFDB_Time t);
\end{lstlisting}
As \code{isigsettime}, for the signals of group \code{sgroup} only (the
signals stored in one file); the other groups keep their position. It serves
programs that open several records at once (\code{isigopen("+REC", ...)}).
\begin{params}
\item[sgroup] The signal group number.
\item[t] The time, in sample intervals.
\end{params}
\returns{0 on success; -1 at the end of the data or on an improper seek; -2
if \code{sgroup} is not a signal group.}
\raises{\code{isigsettime: improper seek on signal group N};
\code{isgsettime: incorrect signal group number N}.}

WFDB-x: a group number beyond the groups of the open record is reported as
incorrect (-2) also when the signals are already at \code{t}, where C reads
past the end of its array (\cref{chap:behavior}).

\apientry{tnextvec}{Find the next valid sample of a signal}
\begin{lstlisting}[language=C,style=signature]
WFDB_Time tnextvec(WFDB_Signal s, WFDB_Time t);
\end{lstlisting}
Finds the first time, at or after \code{t}, at which input signal \code{s}
has a valid sample (not \code{WFDB\_INVALID\_SAMPLE}), and positions the
input signals there, so that the next \code{getvec} reads it. It skips long
gaps (in a multi-segment record, segments without the signal) efficiently.
\begin{params}
\item[s] The input signal number.
\item[t] The time at which the search begins, in sample intervals.
\end{params}
\returns{The time of the next valid sample (\code{>= 0}); -1 at the end of
the data, on an improper seek, or if \code{s} is not an input signal (or the
negative status of \code{getvec}).}
\raises{\code{nextvect: illegal signal number N} (the C library's spelling).}

\apientry{sample}{Read a sample by its time}
\begin{lstlisting}[language=C,style=signature]
WFDB_Sample sample(WFDB_Signal s, WFDB_Time t);
\end{lstlisting}
Returns sample \code{t} of input signal \code{s}, in any order: the samples
are read with \code{getvec} (which moves the input signals) into a buffer of
the 4096 most recent sample vectors, from which they are returned. A negative
\code{t} is taken as 0. \code{sample\_valid} tells whether the value is
valid. \code{wfdbquit} releases the buffer. In R, this function is
\code{wfdb\_sample}.
\begin{params}
\item[s] The input signal number (0 to the number of input signals minus 1).
\item[t] The sample number.
\end{params}
\returns{The sample, in ADC units; after the end of the record, the last
sample read; \code{WFDB\_INVALID\_SAMPLE} if the signal has no valid sample at
\code{t}, if \code{s} is not an input signal, or if the signals cannot be
positioned.}

\apientry{sample\_valid}{Whether the last sample read by sample is valid}
\begin{lstlisting}[language=C,style=signature]
int sample_valid(void);
\end{lstlisting}
Tells whether the value most recently returned by \code{sample} is valid, so
that \code{while (sample\_valid())} reads to the end of the record and
\code{sample\_valid() > 0} selects the values usable in a computation.
\returns{1 if the value was valid; 0 if \code{t} was after the end of the
record (the value is padding); -1 if the value was
\code{WFDB\_INVALID\_SAMPLE} (the signal is not available at \code{t}, or the
signal number is not valid).}

\apientry{findsig}{Find an input signal by name or number}
\begin{lstlisting}[language=C,style=signature]
int findsig(const char *signame);
\end{lstlisting}
Returns the number of the input signal designated by \code{signame}: a string
of decimal digits that is the number of an open input signal is taken as that
number; otherwise the lowest-numbered input signal whose description
(\code{desc}) is exactly \code{signame} is returned. The descriptions are
compared as the header gives them, spaces included (the high-level readers
remove surrounding spaces from signal names; the library does not).
\begin{params}
\item[signame] A signal number or a signal description, such as \code{"MLII"}.
\end{params}
\returns{The signal number, or -1 if there is no such signal.}

\apientry{getspf}{Number of samples per signal per frame}
\begin{lstlisting}[language=C,style=signature]
int getspf(void);
\end{lstlisting}
Returns the number of sample vectors that \code{getvec} returns per frame:
1, except in \code{WFDB\_HIGHRES} mode for a record with oversampled signals,
where it is the largest number of samples per frame of the input signals (the
sampling frequency divided by the frame frequency, rounded).
\returns{The number of samples per signal per frame.}

\apientry{setgvmode}{Set the mode of getvec}
\begin{lstlisting}[language=C,style=signature]
void setgvmode(int mode);
\end{lstlisting}
Sets how \code{getvec} reads records whose signals have different numbers of
samples per frame, and how invalid samples are returned.
\code{WFDB\_LOWRES} (0, the default): one sample per signal per frame, the
oversampled signals being averaged over the frame. \code{WFDB\_HIGHRES}
(1): every sample of the most oversampled signal, the other signals being
repeated (zero-order hold); the sampling frequency returned by
\code{sampfreq(NULL)}, and the times of the time conversion functions and of
the annotations, are then in units of the high-resolution sample interval
(for annotation files opened after the call). \code{WFDB\_GVPAD} (2), added to
either: an invalid or missing sample is replaced by the previous valid sample
of its signal, in \code{getvec}, \code{getframe} and \code{sample}. A negative
\code{mode} restores the default, read from the environment variable
\code{WFDBGVMODE} (or \code{WFDB\_LOWRES} if it is not set).
\begin{params}
\item[mode] \code{WFDB\_LOWRES} or \code{WFDB\_HIGHRES}, optionally plus
  \code{WFDB\_GVPAD}; or a negative value.
\end{params}

\apientry{getgvmode}{The mode of getvec}
\begin{lstlisting}[language=C,style=signature]
int getgvmode(void);
\end{lstlisting}
Returns the mode set by \code{setgvmode}: the \code{WFDB\_HIGHRES} bit is
set in high-resolution mode, the \code{WFDB\_GVPAD} bit if invalid samples
are replaced.
\returns{The mode (0 to 3).}

\apientry{setifreq}{Resample the input signals}
\begin{lstlisting}[language=C,style=signature]
int setifreq(WFDB_Frequency freq);
\end{lstlisting}
Makes \code{getvec} return the input signals resampled at \code{freq} samples
per second, by linear interpolation between the two nearest samples
(\code{getframe} is not affected). It must be called after the signals are
opened. From then on, the time conversion functions, \code{isigsettime} and
\code{isgsettime} use sample intervals at \code{freq}, the annotators opened
afterwards return times at \code{freq}, and the time resolution of the output
annotation files opened afterwards is \code{freq} (it calls
\code{setafreq(freq)}). The resampling uses the greatest common divisor of the
two frequencies, treated as multiples of 0.005 Hz.
\begin{params}
\item[freq] The sampling frequency wanted, in Hz (> 0).
\end{params}
\returns{0 on success; -1 if no record is open, if \code{freq} is not
positive, or (WFDB-x) if the sampling frequency of the open record is not
specified.}
\raises{\code{setifreq: no open input record}; \code{setifreq: improper
frequency F (must be > 0)}; \code{setifreq: sampling frequency is not
specified}.}

WFDB-x: a record without a sampling frequency is an error
(\cref{chap:behavior}); in legacy mode, it is reported as \code{setifreq: no
open input record}, as in C.

\apientry{getifreq}{The frequency of the samples returned by getvec}
\begin{lstlisting}[language=C,style=signature]
WFDB_Frequency getifreq(void);
\end{lstlisting}
Returns the sampling frequency of the samples returned by \code{getvec}: the
one set by \code{setifreq}, or else that of the open record (as
\code{sampfreq(NULL)}).
\returns{The frequency, in Hz (0 if no record is open, or if its frequency is
not specified).}

\apientry{setibsize}{Set the size of the input buffers}
\begin{lstlisting}[language=C,style=signature]
int setibsize(int input_buffer_size);
\end{lstlisting}
Sets the size of the buffers in which the input signal files are read, for
the signals opened afterwards (a block size given in the header takes
precedence). It cannot be used while input signals are open.
\begin{params}
\item[input\_buffer\_size] The size in bytes; 0 for the default
  (\code{BUFSIZ}, 8192 bytes).
\end{params}
\returns{The new size; -1 if input signals are open; -2 if the size is
negative.}
\raises{\code{setibsize: can't change buffer size after isigopen};
\code{setibsize: illegal buffer size N}.}

\apientry{setobsize}{Set the size of the output buffers}
\begin{lstlisting}[language=C,style=signature]
int setobsize(int output_buffer_size);
\end{lstlisting}
As \code{setibsize}, for the output signal files opened afterwards by
\code{osigopen} and \code{osigfopen}; it cannot be used while output signals
are open.
\begin{params}
\item[output\_buffer\_size] The size in bytes; 0 for the default
  (\code{BUFSIZ}, 8192 bytes).
\end{params}
\returns{The new size; -1 if output signals are open; -2 if the size is
negative.}
\raises{\code{setobsize: can't change buffer size after osig[f]open};
\code{setobsize: illegal buffer size N}.}

% ==========================================================================
\section{Time and date conversion}
\label{sec:library-time}

Times are kept as \code{WFDB\_Time} values, in sample intervals. The
functions of this section convert them to and from text, using the sampling
frequency (or the frequency set by \code{setifreq}), the base time and date
and the counter frequency of the open record, set by \code{isigopen},
\code{wfdbinit} or \code{sampfreq} from a header, or by \code{setsampfreq},
\code{setbasetime} and \code{setcfreq}. A positive time is an elapsed time
(an interval from the beginning of the record); a time of day is represented
by the negated number of sample intervals from the beginning of the record
(0 or negative).

WFDB-x (\cref{chap:behavior}): when no sampling frequency is known, C took a
sample interval to be one second; WFDB-x reports \code{FUNCTION: sampling
frequency is not specified}, and \code{timstr} and \code{mstimstr} return the
time as a sample number (\code{s} followed by its magnitude), \code{strtim}
returns 0 for a string that needs the frequency (\code{wfdbx\_time\_needs\_freq}
tells which). Legacy mode keeps C's behavior.

\apientry{timstr}{Format a time}
\begin{lstlisting}[language=C,style=signature]
char *timstr(WFDB_Time t);
\end{lstlisting}
Converts \code{t} to text. If \code{t} is positive, or if neither a base time
nor a base date is defined, \code{|t|} is an elapsed time, written
\code{h:mm:ss} right-aligned in 8 characters (more from 100 hours on),
without the hours under one hour (\code{"    1:00"}, \code{" 1:02:05"}). Otherwise \code{-t} sample intervals
after the base time are written as a time of day in brackets, with all the
digits, followed by the date if a base date is defined, or else by the
number of days after the base time if it is not 0: \code{"[12:00:01
01/01/2020]"}, \code{"[12:00:00 1]"}. The time is rounded to the
millisecond, and the milliseconds are then dropped. A
base time of midnight (0) without a base date counts as undefined.
\begin{params}
\item[t] The time, in sample intervals.
\end{params}
\returns{The text. In C, it is in a static buffer (shared with
\code{mstimstr}) that the next call overwrites; in WFDB-x, it is a new string
in every language.}
\raises{WFDB-x: \code{timstr: sampling frequency is not specified}, and the
result \code{s} followed by \code{|t|} (\code{"s100"}), if the frequency is
not known.}

\apientry{mstimstr}{Format a time with milliseconds}
\begin{lstlisting}[language=C,style=signature]
char *mstimstr(WFDB_Time t);
\end{lstlisting}
As \code{timstr}, with the milliseconds (rounded): an elapsed time
\code{h:mm:ss.mmm} right-aligned in 12 characters (\code{"    0:01.000"}), a
time of day \code{"[hh:mm:ss.mmm dd/mm/yyyy]"}. \code{mstimstr(0)} gives the
base time and date of the record (there is no \code{getbasetime}).
\begin{params}
\item[t] The time, in sample intervals.
\end{params}
\returns{The text, as for \code{timstr}.}
\raises{WFDB-x: \code{mstimstr: sampling frequency is not specified}, as for
\code{timstr}.}

\apientry{strtim}{Parse a time}
\begin{lstlisting}[language=C,style=signature]
WFDB_Time strtim(const char *time_string);
\end{lstlisting}
Converts a time string to sample intervals. Leading spaces, tabs, carriage
returns and newlines are skipped; then, according to the first character:
\begin{description}
\item[\code{s}\emph{N}] sample number \emph{N} (\code{s12345});
\item[\code{c}\emph{X}] counter value \emph{X}, converted with the counter
  frequency and the base counter value (\code{c350.5}); without a counter
  frequency, \emph{X} is a sample number;
\item[\code{f}\emph{N}] frame number \emph{N}, converted to sample intervals;
\item[\code{e}] the end of the record (its length; 0 if the header does not
  give it), in high-resolution intervals in \code{WFDB\_HIGHRES} mode;
\item[\code{i}] the time of the next sample of the input signals;
\item[\code{o}] the time of the next sample of the output signals;
\item[\code{[}\emph{time} \emph{date}\code{]}] a time of day, \code{[hh:mm:ss]},
  followed optionally by a date \code{dd/mm/yyyy} or by a number of days after
  the base date (\code{[13:6:0]}, \code{[8:0:0 1]}, \code{[12:0:0 1/3/1992]}):
  the result is the negated number of sample intervals from the base time to
  that time (0 or negative); without the closing bracket, the result is 0;
\item[otherwise] an elapsed time \code{h:m:s}, \code{m:s} or \code{s}, each
  part with optional decimals (\code{2:14.875}, \code{143}, \code{4:02:01}),
  rounded to the nearest sample interval.
\end{description}
For records with sampling frequencies below 1 kHz,
\code{strtim(mstimstr(t))} is \code{t}.
\begin{params}
\item[time\_string] The time string.
\end{params}
\returns{A positive number of sample intervals for an elapsed time, a sample
number, a counter value, \code{e}, \code{i} or \code{o}; the negated number
of sample intervals for a time of day; 0 for a string that is not a time
(indistinguishable from a time of 0).}
\raises{WFDB-x: \code{strtim: sampling frequency is not specified}, and 0, if
the frequency is not known and the string needs it.}

\apientry{datstr}{Format a date}
\begin{lstlisting}[language=C,style=signature]
char *datstr(WFDB_Date d);
\end{lstlisting}
Converts a Julian date (a day number; these Julian dates begin at midnight
local time, not at noon) to text, \code{" dd/mm/yyyy"}, with a leading space,
which the Programmer's Guide does not mention (\code{datstr(1753639)} is
\code{" 14/07/1789"}). The Gregorian calendar applies from 15 October 1582;
years before 1 A.D. are negative (there is no year 0).
\begin{params}
\item[d] The Julian date.
\end{params}
\returns{The text (a static buffer in C; a new string in WFDB-x).}

\apientry{strdat}{Parse a date}
\begin{lstlisting}[language=C,style=signature]
WFDB_Date strdat(const char *date_string);
\end{lstlisting}
Converts a date \code{dd/mm/yyyy} to a Julian date. The year is taken as
written: \code{15/3/89} is in the year 89, not 1989. The number of days
between two dates is the difference of their Julian dates.
\begin{params}
\item[date\_string] The date.
\end{params}
\returns{The Julian date; 0 if the string does not contain two \code{/}, the
day is not 1 to 31, the month is not 1 to 12, or the year is 0.}

\apientry{sampfreq}{Sampling frequency of a record}
\begin{lstlisting}[language=C,style=signature]
WFDB_Frequency sampfreq(char *record);
\end{lstlisting}
With a record name, reads the header of \code{record} and returns its
sampling frequency; the frequency, base time and date and counter frequency
of the header then become those used by the time conversion functions (as
after \code{isigopen(record, NULL, 0)}, without closing the open signals).
With \code{NULL}, returns the sampling frequency in use. For a record with
oversampled signals, the result is the frame frequency in
\code{WFDB\_LOWRES} mode and the highest sampling frequency in
\code{WFDB\_HIGHRES} mode.
\begin{params}
\item[record] The record name, or \code{NULL}.
\end{params}
\returns{The sampling frequency in Hz; -1 if the header cannot be read; -2 if
it is not valid; WFDB-x: 0 if the header does not give a frequency (C: 250,
\code{WFDB\_DEFFREQ}; see \cref{chap:behavior}). \code{sampfreq(record)}
first sets the frequency to 0, so that a frequency set by \code{setsampfreq}
does not apply to it.}

\apientry{setsampfreq}{Set the sampling frequency}
\begin{lstlisting}[language=C,style=signature]
int setsampfreq(WFDB_Frequency sampling_frequency);
\end{lstlisting}
Sets the frame frequency used by the time conversion functions and written
by \code{newheader} and \code{setheader}: it is the sampling frequency of a
record whose signals have one sample per frame (in \code{WFDB\_HIGHRES} mode,
the sampling frequency becomes the frame frequency times the largest number
of samples per frame). A header read afterwards by \code{isigopen}, \code{osigopen} or
\code{wfdbinit} replaces it, unless the
header gives no frequency: \code{setsampfreq(250)} before opening a record
whose header gives none reproduces C's default (\cref{chap:behavior}).
\begin{params}
\item[sampling\_frequency] The frame frequency, in Hz (0 or more).
\end{params}
\returns{0 on success; -1 if the frequency is negative.}
\raises{\code{setsampfreq: sampling frequency must not be negative}.}

\apientry{setbasetime}{Set the base time and date}
\begin{lstlisting}[language=C,style=signature]
int setbasetime(char *time_string);
\end{lstlisting}
Sets the base time (the time of day of sample 0) and the base date, used by
the time conversion functions and written by \code{newheader} and
\code{setheader}. The string is a time \code{hh:mm:ss} (with optional
decimals), optionally followed by a space and a date \code{dd/mm/yyyy};
leading spaces are skipped. Only a space separates the date (not a tab, which
the Programmer's Guide also allows). An empty string or \code{NULL} sets the
current local time and date.
\begin{params}
\item[time\_string] The time and date, or \code{""} or \code{NULL}.
\end{params}
\returns{0 on success; -1 if neither a time nor a date can be read from the
string (and it does not begin with \code{[}).}
\raises{\code{setbasetime: incorrect time format, 'STRING'}.}

\apientry{getcfreq}{Counter frequency}
\begin{lstlisting}[language=C,style=signature]
WFDB_Frequency getcfreq(void);
\end{lstlisting}
Returns the counter frequency: the rate at which a counter (of a tape
recorder, or the pages of a chart) changes, as given by the header or by
\code{setcfreq}, which \code{strtim} uses for \code{c} times. If none is
defined (0 or negative), it returns the frame frequency.
\returns{The counter frequency, in Hz.}

\apientry{setcfreq}{Set the counter frequency}
\begin{lstlisting}[language=C,style=signature]
void setcfreq(WFDB_Frequency counter_frequency);
\end{lstlisting}
Sets the counter frequency, used by \code{strtim} and written by
\code{newheader} and \code{setheader}; 0 or a negative value means none. A
header read afterwards replaces it. \code{wfdbquit} does not reset it.
\begin{params}
\item[counter\_frequency] The counter frequency, in Hz.
\end{params}

\apientry{getbasecount}{Base counter value}
\begin{lstlisting}[language=C,style=signature]
double getbasecount(void);
\end{lstlisting}
Returns the counter value of sample 0, as given by the header or by
\code{setbasecount} (0 if neither gave one).
\returns{The base counter value.}

\apientry{setbasecount}{Set the base counter value}
\begin{lstlisting}[language=C,style=signature]
void setbasecount(double count);
\end{lstlisting}
Sets the counter value of sample 0, used by \code{strtim} and written by
\code{newheader} and \code{setheader}. A header read afterwards replaces it.
\code{wfdbquit} does not reset it.
\begin{params}
\item[count] The base counter value.
\end{params}

% ==========================================================================
\section{Amplitude conversion}
\label{sec:library-amplitude}

These functions convert between ADC units and physical units with the gain
(and, for \code{aduphys} and \code{physadu}, the baseline) of an open input
signal. \code{adumuv} and \code{muvadu} convert differences of level (the
baseline is not used), between ADC units and thousandths of the physical
unit (microvolts for a signal in millivolts).

WFDB-x (\cref{chap:behavior}): for an uncalibrated signal (gain 0) or a
number that is not that of an open input signal, C assumed a gain of 200
(\code{WFDB\_DEFGAIN}) and a baseline of 0; WFDB-x reports \code{FUNCTION:
signal N is uncalibrated} or \code{FUNCTION: invalid signal number N}, and
returns \code{WFDB\_INVALID\_SAMPLE} (\code{aduphys}: NaN). Legacy mode keeps
C's behavior.

\apientry{aduphys}{ADC units to physical units}
\begin{lstlisting}[language=C,style=signature]
double aduphys(WFDB_Signal s, WFDB_Sample a);
\end{lstlisting}
Returns \code{(a - baseline) / gain} for input signal \code{s}.
\begin{params}
\item[s] The input signal number.
\item[a] The sample value, in ADC units.
\end{params}
\returns{The value in physical units; WFDB-x: NaN on error.}

\apientry{physadu}{Physical units to ADC units}
\begin{lstlisting}[language=C,style=signature]
WFDB_Sample physadu(WFDB_Signal s, double v);
\end{lstlisting}
Returns \code{v * gain}, rounded to the nearest integer (halves away from
zero), plus the baseline of input signal \code{s}.
\begin{params}
\item[s] The input signal number.
\item[v] The value, in physical units.
\end{params}
\returns{The sample value in ADC units; WFDB-x: \code{WFDB\_INVALID\_SAMPLE}
on error.}

\apientry{adumuv}{ADC units to microvolts}
\begin{lstlisting}[language=C,style=signature]
int adumuv(WFDB_Signal s, WFDB_Sample a);
\end{lstlisting}
Returns \code{a * 1000 / gain} for input signal \code{s}, rounded to the
nearest integer (halves away from zero).
\begin{params}
\item[s] The input signal number.
\item[a] The difference of level, in ADC units.
\end{params}
\returns{The difference in thousandths of the physical unit; WFDB-x:
\code{WFDB\_INVALID\_SAMPLE} on error.}

\apientry{muvadu}{Microvolts to ADC units}
\begin{lstlisting}[language=C,style=signature]
WFDB_Sample muvadu(WFDB_Signal s, int microvolts);
\end{lstlisting}
Returns \code{microvolts * gain / 1000} for input signal \code{s}, rounded to
the nearest integer (halves away from zero).
\begin{params}
\item[s] The input signal number.
\item[microvolts] The difference of level, in thousandths of the physical
  unit.
\end{params}
\returns{The difference in ADC units; WFDB-x: \code{WFDB\_INVALID\_SAMPLE} on
error.}

% ==========================================================================
\section{Annotations}
\label{sec:library-annotations}

Annotation files are opened by \code{annopen} (or \code{wfdbinit}), read
sequentially with \code{getann}, and written with \code{putann}, in MIT
format (\code{WFDB\_READ}, \code{WFDB\_WRITE}) or AHA format
(\code{WFDB\_AHA\_READ}, \code{WFDB\_AHA\_WRITE}); \cref{chap:formats}
describes both formats and lists the annotation codes. An annotation is a
\code{WFDB\_Annotation} structure (\cref{sec:library-structures}).

\apientry{getann}{Read the next annotation}
\begin{lstlisting}[language=C,style=signature]
int getann(WFDB_Annotator a, WFDB_Annotation *annot);
\end{lstlisting}
Reads the next annotation of input annotator \code{a} into \code{annot} (an
annotation pushed back by \code{ungetann} first, if there is one). Its time
is in sample intervals of the input sampling frequency in effect when the
annotator was opened (that of \code{setifreq}, or of the record, in
high-resolution intervals in \code{WFDB\_HIGHRES} mode), converted from the
time resolution recorded in the file if it records one, or in the units set by
\code{setiafreq}. In C, \code{aux} points to a buffer of the annotator that
later calls overwrite.
\begin{params}
\item[a] The input annotator number (0 for the first input annotator of
  \code{annopen}).
\item[annot] The annotation structure, filled by the function.
\end{params}
\returns{0 on success; -1 at the end of the file (\code{*annot} is not
valid); -2 if \code{a} is not an open input annotator; -3 at an unexpected end
of the file.}
\raises{\code{getann: can't read annotator N}; \code{getann: unexpected EOF
in annotator ATR}; \code{getann: unexpected annot number in annotator ATR}.}

\apientry{ungetann}{Push an annotation back}
\begin{lstlisting}[language=C,style=signature]
int ungetann(WFDB_Annotator a, const WFDB_Annotation *annot);
\end{lstlisting}
Makes \code{*annot} the next annotation that \code{getann} returns for input
annotator \code{a}. It need not have been read from that annotator. Only one
annotation can be pushed back at a time for each annotator.
\begin{params}
\item[a] The input annotator number.
\item[annot] The annotation to push back.
\end{params}
\returns{0 on success; -1 if an annotation is already pushed back (this one
is not); -2 if \code{a} is not an input annotator.}
\raises{\code{ungetann: annotator N is not initialized}; \code{ungetann:
pushback buffer is full}.}

\apientry{putann}{Write an annotation}
\begin{lstlisting}[language=C,style=signature]
int putann(WFDB_Annotator a, const WFDB_Annotation *annot);
\end{lstlisting}
Writes \code{*annot} to output annotator \code{a}; all its fields must be
set. Annotations may be written in any order: an annotation is in order if
its time is later than that of the previous annotation, or its \code{chan} or
its \code{num} is greater, or both are at time 0; a file that received one out
of order is sorted when it is closed (see \code{oannclose}). The first
\code{putann} to a file writes, before the annotation, the time resolution set
by \code{setafreq} and the modification labels of the mnemonics and
descriptions changed by \code{setannstr} and \code{setanndesc}. In an AHA
file, at most 6 bytes of \code{aux} are written.
\begin{params}
\item[a] The output annotator number (0 for the first output annotator of
  \code{annopen}).
\item[annot] The annotation, filled by the caller.
\end{params}
\returns{0 on success; -1 on a write error, or for a time equal to
\code{WFDB\_TIME\_MIN} or \code{WFDB\_TIME\_MAX}; -2 if \code{a} is not an
open output annotator.}
\raises{\code{putann: can't write annotation file N}; \code{putann: write
error on annotation file ATR}; \code{putann: time overflow in annotation file
N}.}

\apientry{iannsettime}{Skip to a time in all input annotators}
\begin{lstlisting}[language=C,style=signature]
int iannsettime(WFDB_Time t);
\end{lstlisting}
Positions every input annotator at its first annotation at or after time
\code{|t|} (rewinding the file if needed, which a pipe cannot do), so that a
value of \code{strtim} for a time of day can be given as it is.
\code{WFDB\_TIME\_MIN} rewinds to the very beginning, before annotations at
negative times. A pushed-back annotation is discarded if the file is rewound.
\begin{params}
\item[t] The time, in the units of the annotators' times.
\end{params}
\returns{0 if at least one annotator is positioned on an annotation; -1 if
every annotator reached the end of its file (or none is open), or on an
improper seek. (The Programmer's Guide also lists -3, which the C code never
returns.)}
\raises{\code{iannsettime: improper seek}.}

\apientry{iannclose}{Close an input annotator}
\begin{lstlisting}[language=C,style=signature]
void iannclose(WFDB_Annotator a);
\end{lstlisting}
Closes the file of input annotator \code{a}. The input annotators that follow
are renumbered: annotator \code{a+1} becomes \code{a}, and so on. Nothing is
done if \code{a} is not open.
\begin{params}
\item[a] The input annotator number.
\end{params}

\apientry{oannclose}{Close an output annotator}
\begin{lstlisting}[language=C,style=signature]
void oannclose(WFDB_Annotator a);
\end{lstlisting}
Completes and closes the file of output annotator \code{a} (an end marker in
an MIT file, padding to a 1024-byte block in an AHA file); the output
annotators that follow are renumbered. If annotations were written out of
order, the file is sorted, unless the environment variable \code{WFDBANNSORT}
is 0: the result is that of \code{sortann -r. REC -a ATR}, with the messages
\code{Rearranging annotations for output annotator ATR ...done!}; if it is
not sorted, the advice \code{Use the command: sortann -r REC -a ATR} is
printed. The status is kept for \code{annopen(..., 0)}: -7 after a write
error, -6 after an out-of-order file.
\begin{params}
\item[a] The output annotator number.
\end{params}
\raises{\code{oannclose: write error on annotation file ATR};
\code{Annotations still need to be rearranged.}}

WFDB-x (\cref{chap:behavior}): the file is sorted in the same process (C ran
the program \code{sortann} through the shell), and an annotation at time 0,
number 0 and channel 0 that arrives after later ones is no longer lost, nor
are those written after it; legacy mode runs the \code{sortann} tool of the
same language, with C's result (without a shell; R and MATLAB/Octave quote
its arguments instead). A write error is counted
(\code{wfdbx\_write\_errors}). The high-level \code{write\_annotations} can
decide for its call whether the file is sorted.

\apientry{setafreq}{Set the time resolution of output annotation files}
\begin{lstlisting}[language=C,style=signature]
void setafreq(WFDB_Frequency f);
\end{lstlisting}
Sets the time resolution, in ticks per second, of the output annotation files
to which no annotation has been written yet (contrary to the Programmer's
Guide, this includes those already open): the resolution is recorded at the
beginning of the file, and the times given to \code{putann} are in ticks. 0
(the default) records none: the times are then in sample intervals of the
record. \code{setifreq} calls it. \code{wfdbquit} does not reset it.
\begin{params}
\item[f] The time resolution, in Hz; 0 for none.
\end{params}

\apientry{getafreq}{Time resolution of output annotation files}
\begin{lstlisting}[language=C,style=signature]
WFDB_Frequency getafreq(void);
\end{lstlisting}
Returns the value set by \code{setafreq} (or by \code{setifreq}).
\returns{The time resolution in Hz; 0 if none was set.}

\apientry{setiafreq}{Set the time units of an input annotator}
\begin{lstlisting}[language=C,style=signature]
void setiafreq(WFDB_Annotator a, WFDB_Frequency f);
\end{lstlisting}
Makes \code{getann} return the times of input annotator \code{a} in units of
\code{1/f} second (the times read and pushed back so far are converted too).
\code{annopen} sets it to the input sampling frequency. With \code{f} 0 or
negative, the times are those of the file, unconverted. Nothing is done if
\code{a} is not an open input annotator.
\begin{params}
\item[a] The input annotator number.
\item[f] The frequency, in Hz.
\end{params}

\apientry{getiafreq}{Time units of an input annotator}
\begin{lstlisting}[language=C,style=signature]
WFDB_Frequency getiafreq(WFDB_Annotator a);
\end{lstlisting}
Returns the frequency of the time units that \code{getann} uses for input
annotator \code{a} (see \code{setiafreq}).
\returns{The frequency in Hz; -2 if \code{a} is not an open input annotator.}

\apientry{getiaorigfreq}{Time resolution recorded in an annotation file}
\begin{lstlisting}[language=C,style=signature]
WFDB_Frequency getiaorigfreq(WFDB_Annotator a);
\end{lstlisting}
Returns the time resolution recorded in the file of input annotator \code{a}
by the program that wrote it (see \code{setafreq}).
\returns{The time resolution in Hz; 0 if the file records none (its times are
then in sample intervals of the record's frames); -2 if \code{a} is not an
open input annotator.}

\apientry{annstr}{Mnemonic of an annotation code}
\begin{lstlisting}[language=C,style=signature]
char *annstr(int annotation_code);
\end{lstlisting}
Returns the mnemonic of an annotation code (\code{"N"} for \code{NORMAL}),
from the table that \code{setannstr} and the modification labels of the input
annotation files change (\cref{chap:formats} lists the standard mnemonics).
The codes of the range 0 to \code{ACMAX} that are not defined have mnemonics
such as \code{"[15]"}.
\begin{params}
\item[annotation\_code] The annotation code.
\end{params}
\returns{The mnemonic; for a code outside 0 to \code{ACMAX}, the code in
brackets (\code{"[55]"}).}

\apientry{strann}{Annotation code of a mnemonic}
\begin{lstlisting}[language=C,style=signature]
int strann(const char *annotation_mnemonic_string);
\end{lstlisting}
Returns the code (1 to \code{ACMAX}) whose mnemonic, in the table of
\code{annstr}, is \code{annotation\_mnemonic\_string}; if several codes have
it, the lowest.
\begin{params}
\item[annotation\_mnemonic\_string] The mnemonic.
\end{params}
\returns{The annotation code; \code{NOTQRS} (0) if no code has that mnemonic.}

\apientry{setannstr}{Change the mnemonic of an annotation code}
\begin{lstlisting}[language=C,style=signature]
int setannstr(int annotation_code, const char *annotation_mnemonic_string);
\end{lstlisting}
Sets the mnemonic of a code in the table of \code{annstr} and \code{strann}
(\code{NULL} is taken as \code{""}). With a positive code, the change is also
marked to be recorded, as a modification label, in the output annotation
files opened afterwards, so that other programs read the files with these
mnemonics. With a negative code, the mnemonic of \code{-code} is changed for
this program only. \code{wfdbquit} does not reset the table.
\begin{params}
\item[annotation\_code] The code (or its negation).
\item[annotation\_mnemonic\_string] The new mnemonic.
\end{params}
\returns{0 on success; -1 if \code{|code|} is greater than \code{ACMAX}.}
\raises{\code{setannstr: illegal annotation code N}.}

\apientry{anndesc}{Description of an annotation code}
\begin{lstlisting}[language=C,style=signature]
char *anndesc(int annotation_code);
\end{lstlisting}
Returns the description of an annotation code (\code{"Normal beat"} for
\code{NORMAL}), from the table that \code{setanndesc} and the modification
labels of the input annotation files change.
\begin{params}
\item[annotation\_code] The annotation code.
\end{params}
\returns{The description; \code{NULL} for a code of the range 0 to
\code{ACMAX} that has none; \code{"illegal annotation code"} for a code outside
that range (the Programmer's Guide says \code{NULL}; the C code returns this
string).}

\apientry{setanndesc}{Change the description of an annotation code}
\begin{lstlisting}[language=C,style=signature]
int setanndesc(int annotation_code, const char *annotation_description);
\end{lstlisting}
As \code{setannstr}, for the table of \code{anndesc}: a positive code also
marks the description to be recorded in the output annotation files opened
afterwards, a negative code changes the description of \code{-code} for this
program only.
\begin{params}
\item[annotation\_code] The code (or its negation).
\item[annotation\_description] The new description.
\end{params}
\returns{0 on success; -1 if \code{|code|} is greater than \code{ACMAX}.}
\raises{\code{setanndesc: illegal annotation code N}.}

\apientry{ecgstr}{Standard mnemonic of an annotation code}
\begin{lstlisting}[language=C,style=signature]
char *ecgstr(int annotation_code);
\end{lstlisting}
As \code{annstr}, from a separate table that only \code{setecgstr} changes
(not the modification labels of annotation files), for programs that need a
fixed set of mnemonics.
\begin{params}
\item[annotation\_code] The annotation code.
\end{params}
\returns{The mnemonic; the code in brackets for a code outside 0 to
\code{ACMAX}.}

\apientry{strecg}{Annotation code of a standard mnemonic}
\begin{lstlisting}[language=C,style=signature]
int strecg(const char *annotation_mnemonic_string);
\end{lstlisting}
As \code{strann}, with the table of \code{ecgstr}.
\begin{params}
\item[annotation\_mnemonic\_string] The mnemonic.
\end{params}
\returns{The annotation code; \code{NOTQRS} (0) if no code has that mnemonic.}

\apientry{setecgstr}{Change a standard mnemonic}
\begin{lstlisting}[language=C,style=signature]
int setecgstr(int annotation_code, const char *annotation_mnemonic_string);
\end{lstlisting}
Sets the mnemonic of a code in the table of \code{ecgstr} and \code{strecg}
(\code{NULL} is taken as \code{""}). Nothing is recorded in annotation files,
and negative codes are illegal.
\begin{params}
\item[annotation\_code] The code, 0 to \code{ACMAX}.
\item[annotation\_mnemonic\_string] The new mnemonic.
\end{params}
\returns{0 on success; -1 if the code is outside 0 to \code{ACMAX}.}
\raises{\code{setecgstr: illegal annotation code N}.}

\subsection{Classification and mapping of annotation codes}
\label{sec:library-ecgmap}

\file{<wfdb/ecgmap.h>} defines macros that classify annotation codes and map
them to smaller sets, with tables that a program can change; the C library
provides each of them also as a function, whose name begins with
\code{wfdb\_} (\code{wfdb\_isqrs} for \code{isqrs}). WFDB-x provides both
names in every language, with the same results; the macro names are
functions too (in C++, inline functions of \file{<wfdbx/ecgmap.hpp>}). A
setter stores the new value as a C \code{char}, and returns it as stored:
\code{setisqrs(2, 300)} returns 44. An argument that is not a legal code
(1 to \code{ACMAX}) changes nothing. In R, these functions accept vectors of
codes. \code{isann} (and \code{wfdb\_isann}) returns a logical value in
Python, R and MATLAB/Octave (\code{True}, \code{TRUE}, \code{true}), and 1 or
0 elsewhere.

\apientry{isann}{Whether a code is a legal annotation code}
\begin{lstlisting}[language=C,style=signature]
int isann(int code);              /* a macro of <wfdb/ecgmap.h> */
\end{lstlisting}
Tells whether \code{code} is a legal annotation code, 1 to \code{ACMAX} (49).
\begin{params}
\item[code] The annotation code.
\end{params}
\returns{1 if it is legal, 0 otherwise.}

\apientry{wfdb\_isann}{Function form of isann}
\begin{lstlisting}[language=C,style=signature]
int wfdb_isann(int code);
\end{lstlisting}
The function form of \code{isann}, with the same result.
\begin{params}
\item[code] The annotation code.
\end{params}
\returns{1 if it is legal, 0 otherwise.}

\apientry{isqrs}{Whether a code is a beat}
\begin{lstlisting}[language=C,style=signature]
int isqrs(int code);              /* a macro of <wfdb/ecgmap.h> */
\end{lstlisting}
Tells whether annotation code \code{code} denotes a QRS complex (a beat), from
a table that \code{setisqrs} changes. Initially, the beat codes are 1 to 13,
25, 30, 31, 34, 35, 38 and 41.
\begin{params}
\item[code] The annotation code.
\end{params}
\returns{The value of the table (initially 1 for a beat, 0 otherwise); 0 for
an illegal code.}

\apientry{wfdb\_isqrs}{Function form of isqrs}
\begin{lstlisting}[language=C,style=signature]
int wfdb_isqrs(int code);
\end{lstlisting}
The function form of \code{isqrs}, with the same result.
\begin{params}
\item[code] The annotation code.
\end{params}
\returns{As \code{isqrs}.}

\apientry{setisqrs}{Change the result of isqrs}
\begin{lstlisting}[language=C,style=signature]
int setisqrs(int code, int newval);   /* a macro of <wfdb/ecgmap.h> */
\end{lstlisting}
Sets the value that \code{isqrs(code)} returns.
\begin{params}
\item[code] The annotation code.
\item[newval] The new value.
\end{params}
\returns{The value stored (\code{newval} as a C \code{char}); 0 for an
illegal code.}

\apientry{wfdb\_setisqrs}{Function form of setisqrs}
\begin{lstlisting}[language=C,style=signature]
int wfdb_setisqrs(int code, int newval);
\end{lstlisting}
The function form of \code{setisqrs}, with the same result.
\begin{params}
\item[code] The annotation code.
\item[newval] The new value.
\end{params}
\returns{As \code{setisqrs}.}

\apientry{map1}{Map a code to the five basic classes}
\begin{lstlisting}[language=C,style=signature]
int map1(int code);               /* a macro of <wfdb/ecgmap.h> */
\end{lstlisting}
Maps an annotation code to one of \code{NOTQRS}, \code{NORMAL}, \code{PVC},
\code{FUSION} and \code{LEARN}, from a table that \code{setmap1} changes (for
example, the bundle branch block beats map to \code{NORMAL}, ventricular
escape beats to \code{PVC}).
\begin{params}
\item[code] The annotation code.
\end{params}
\returns{The mapped code; \code{NOTQRS} for an illegal code.}

\apientry{wfdb\_map1}{Function form of map1}
\begin{lstlisting}[language=C,style=signature]
int wfdb_map1(int code);
\end{lstlisting}
The function form of \code{map1}, with the same result.
\begin{params}
\item[code] The annotation code.
\end{params}
\returns{As \code{map1}.}

\apientry{setmap1}{Change the result of map1}
\begin{lstlisting}[language=C,style=signature]
int setmap1(int code, int newval);    /* a macro of <wfdb/ecgmap.h> */
\end{lstlisting}
Sets the value that \code{map1(code)} returns.
\begin{params}
\item[code] The annotation code.
\item[newval] The new value.
\end{params}
\returns{The value stored (as a C \code{char}); \code{NOTQRS} for an illegal
code.}

\apientry{wfdb\_setmap1}{Function form of setmap1}
\begin{lstlisting}[language=C,style=signature]
int wfdb_setmap1(int code, int newval);
\end{lstlisting}
The function form of \code{setmap1}, with the same result.
\begin{params}
\item[code] The annotation code.
\item[newval] The new value.
\end{params}
\returns{As \code{setmap1}.}

\apientry{map2}{Map a code to the six basic classes}
\begin{lstlisting}[language=C,style=signature]
int map2(int code);               /* a macro of <wfdb/ecgmap.h> */
\end{lstlisting}
As \code{map1}, with \code{SVPB} as a sixth class: the supraventricular
premature beats (\code{ABERR}, \code{NPC}, \code{APC}, \code{SVPB}) map to
\code{SVPB} instead of \code{NORMAL}. \code{setmap2} changes its table.
\begin{params}
\item[code] The annotation code.
\end{params}
\returns{The mapped code; \code{NOTQRS} for an illegal code.}

\apientry{wfdb\_map2}{Function form of map2}
\begin{lstlisting}[language=C,style=signature]
int wfdb_map2(int code);
\end{lstlisting}
The function form of \code{map2}, with the same result.
\begin{params}
\item[code] The annotation code.
\end{params}
\returns{As \code{map2}.}

\apientry{setmap2}{Change the result of map2}
\begin{lstlisting}[language=C,style=signature]
int setmap2(int code, int newval);    /* a macro of <wfdb/ecgmap.h> */
\end{lstlisting}
Sets the value that \code{map2(code)} returns.
\begin{params}
\item[code] The annotation code.
\item[newval] The new value.
\end{params}
\returns{The value stored (as a C \code{char}); \code{NOTQRS} for an illegal
code.}

\apientry{wfdb\_setmap2}{Function form of setmap2}
\begin{lstlisting}[language=C,style=signature]
int wfdb_setmap2(int code, int newval);
\end{lstlisting}
The function form of \code{setmap2}, with the same result.
\begin{params}
\item[code] The annotation code.
\item[newval] The new value.
\end{params}
\returns{As \code{setmap2}.}

\apientry{ammap}{Map an AHA annotation code to an MIT code}
\begin{lstlisting}[language=C,style=signature]
int ammap(int code);              /* a macro of <wfdb/ecgmap.h> */
\end{lstlisting}
Maps an AHA annotation code, a character (its character code, such as that of
\code{'V'}), to an MIT annotation code: \code{'N'} to \code{NORMAL},
\code{'V'} to \code{PVC}, \code{'F'} to \code{FUSION}, \code{'E'} to
\code{VESC}, \code{'P'} to \code{PACE}, \code{'Q'} to \code{UNKNOWN},
\code{'R'} to \code{RONT}, \code{'O'} to \code{NOTE}, \code{'U'} to
\code{NOISE}, \code{'['} to \code{VFON}, \code{']'} to \code{VFOFF}. The table
is fixed.
\begin{params}
\item[code] The AHA code.
\end{params}
\returns{The MIT code; \code{NOTQRS} for any other character.}

\apientry{wfdb\_ammap}{Function form of ammap}
\begin{lstlisting}[language=C,style=signature]
int wfdb_ammap(int code);
\end{lstlisting}
The function form of \code{ammap}, with the same result.
\begin{params}
\item[code] The AHA code.
\end{params}
\returns{As \code{ammap}.}

\apientry{mamap}{Map an MIT annotation code to an AHA code}
\begin{lstlisting}[language=C,style=signature]
int mamap(int code, int subtype);     /* a macro of <wfdb/ecgmap.h> */
\end{lstlisting}
Maps an MIT annotation code, with its subtype, to an AHA code (a character
code), from a fixed table: \code{NOISE} maps to \code{'U'} only with subtype
-1 (unreadable signals), and otherwise to \code{'O'}, as do the codes without
an AHA equivalent.
\begin{params}
\item[code] The MIT annotation code.
\item[subtype] Its subtype.
\end{params}
\returns{The AHA code; \code{'O'} for an illegal code.}

\apientry{wfdb\_mamap}{Function form of mamap}
\begin{lstlisting}[language=C,style=signature]
int wfdb_mamap(int code, int subtype);
\end{lstlisting}
The function form of \code{mamap}, with the same result.
\begin{params}
\item[code] The MIT annotation code.
\item[subtype] Its subtype.
\end{params}
\returns{As \code{mamap}.}

\apientry{annpos}{Display position of an annotation code}
\begin{lstlisting}[language=C,style=signature]
int annpos(int code);             /* a macro of <wfdb/ecgmap.h> */
\end{lstlisting}
Returns where programs that display annotations place those of code
\code{code}, from a table that \code{setannpos} changes: \code{APSTD},
\code{APHIGH}, \code{APLOW}, \code{APATT}, \code{APAHIGH}, \code{APALOW} or
\code{APUNDEF} (\cref{tab:library-constants}).
\begin{params}
\item[code] The annotation code.
\end{params}
\returns{The position code; \code{APUNDEF} for an illegal code.}

\apientry{wfdb\_annpos}{Function form of annpos}
\begin{lstlisting}[language=C,style=signature]
int wfdb_annpos(int code);
\end{lstlisting}
The function form of \code{annpos}, with the same result.
\begin{params}
\item[code] The annotation code.
\end{params}
\returns{As \code{annpos}.}

\apientry{setannpos}{Change the result of annpos}
\begin{lstlisting}[language=C,style=signature]
int setannpos(int code, int newval);  /* a macro of <wfdb/ecgmap.h> */
\end{lstlisting}
Sets the value that \code{annpos(code)} returns.
\begin{params}
\item[code] The annotation code.
\item[newval] The new position code.
\end{params}
\returns{The value stored (as a C \code{char}); \code{APUNDEF} for an
illegal code.}

\apientry{wfdb\_setannpos}{Function form of setannpos}
\begin{lstlisting}[language=C,style=signature]
int wfdb_setannpos(int code, int newval);
\end{lstlisting}
The function form of \code{setannpos}, with the same result.
\begin{params}
\item[code] The annotation code.
\item[newval] The new position code.
\end{params}
\returns{As \code{setannpos}.}

% ==========================================================================
\section{Calibration}
\label{sec:library-calibration}

The calibration functions manage the \emph{calibration list}, kept in
memory: for types of signals, the levels and shape of their calibration
pulses and their customary plotting scale (\code{WFDB\_Calinfo},
\cref{sec:library-structures}). A calibration file has one entry per line:
the signal type (\code{*} for any), a tab, then the low and high levels
(\code{-} for an AC-coupled signal's low level, or for a high level equal to
the low level), the pulse shape (\code{square}, \code{sine},
\code{sawtooth}, or anything else for an undefined shape), the scale and the
units, separated by spaces or tabs; lines that begin with \code{\#} and
incomplete lines are ignored.

\apientry{calopen}{Read a calibration file}
\begin{lstlisting}[language=C,style=signature]
int calopen(const char *calibration_filename);
\end{lstlisting}
Reads a calibration file into the calibration list, which is emptied first,
unless the name begins with \code{+}, which is then removed (the entries are
appended). \code{NULL} reads the file named by the environment variable
\code{WFDBCAL}, or else \file{wfdbcal}; \code{-} reads the standard input.
\begin{params}
\item[calibration\_filename] The name of the calibration file, or
  \code{NULL}.
\end{params}
\returns{0 on success; -2 if the file cannot be opened. (The Programmer's
Guide also lists -1, for a lack of memory.)}
\raises{\code{calopen: can't read calibration file FILE}.}

WFDB-x: the file is read only where its name says, under the input directory
for a relative name, not searched for in the WFDB path (\cref{chap:behavior}).

\apientry{getcal}{Find the calibration of a signal}
\begin{lstlisting}[language=C,style=signature]
int getcal(const char *description, const char *units, WFDB_Calinfo *cal);
\end{lstlisting}
Copies into \code{*cal} the first entry of the calibration list whose signal
type is a prefix of \code{description} (or equal to it), and whose units are
\code{units}. A \code{NULL} argument matches every entry.
\begin{params}
\item[description] The signal description (\code{desc} of a
  \code{WFDB\_Siginfo}), or \code{NULL}.
\item[units] The physical units, or \code{NULL}.
\item[cal] The structure filled by the function.
\end{params}
\returns{0 if an entry matches; -1 otherwise (also if the list is empty).}

In R and MATLAB/Octave, \code{getcal} takes \code{description} and
\code{units} only, and returns the entry with the status: \code{list(stat,
cal)} in R, \code{[stat, cal]} in MATLAB/Octave.

\apientry{putcal}{Add an entry to the calibration list}
\begin{lstlisting}[language=C,style=signature]
int putcal(const WFDB_Calinfo *cal);
\end{lstlisting}
Appends a copy of \code{*cal} to the end of the calibration list.
\begin{params}
\item[cal] The entry.
\end{params}
\returns{0. (The Programmer's Guide also lists -1, for a lack of memory.)}

\apientry{newcal}{Write the calibration list to a file}
\begin{lstlisting}[language=C,style=signature]
int newcal(const char *calibration_filename);
\end{lstlisting}
Writes the calibration list, which is not changed, to a new file in the
format read by \code{calopen}, with \code{undefined} for an undefined pulse
shape. The name may contain only letters, digits, \code{\_},
\code{\textasciitilde}, \code{-} and \code{/}.
\begin{params}
\item[calibration\_filename] The name of the file.
\end{params}
\returns{0 on success; -1 if the name is illegal, the file cannot be created,
or it cannot be written completely.}
\raises{\code{newcal: can't create calibration file FILE}; \code{newcal: write
error in calibration file}; \code{init: illegal character N in calibration
file name}.}

WFDB-x: the file is created under the output directory; a write error is
counted (\code{wfdbx\_write\_errors}).

\apientry{flushcal}{Empty the calibration list}
\begin{lstlisting}[language=C,style=signature]
void flushcal(void);
\end{lstlisting}
Empties the calibration list. \code{wfdbquit} does not.

% ==========================================================================
\section{Header and info strings}
\label{sec:library-header}

A header (\file{.hea} file) describes a record: \code{newheader} writes it
for the signals written with \code{putvec}, \code{setheader} from signal
specifications, \code{setmsheader} for a multi-segment record. The header
being written stays open, so that \code{putinfo} can add info strings to it,
until another header is written or \code{osigfopen(NULL, 0)} (or
\code{wfdbquit}) closes it. Info strings are the comment lines of a header
that follow the record line, and the lines of an optional \file{.info} file;
\cref{chap:metadata} describes the \code{Key: value} entries that WFDB-x
reads and writes in them.

WFDB-x (\cref{chap:behavior}): headers are created under the output
directory; a header written under a path (\code{out/rec}) names the record,
its signal files and its segments relative to its own directory
(\code{rec}, \code{rec.dat}), and only the record's own name (the last
component) must be made of letters, digits, \code{\_}, \code{-} and
\code{\textasciitilde}. Legacy mode checks the whole name and writes the
names as given.

\apientry{newheader}{Write the header of the output signals}
\begin{lstlisting}[language=C,style=signature]
int newheader(char *record);
\end{lstlisting}
Writes the header of \code{record} for the open output signals, as
\code{setheader} does, with the numbers of samples and the checksums that
\code{putvec} computed: it is called after the samples are written, before the
output signals are closed. The frequencies and the base time are those set
by \code{setsampfreq}, \code{setcfreq}, \code{setbasecount} and
\code{setbasetime}, and the skews and prolog lengths those set by
\code{wfdbsetskew} and \code{wfdbsetstart}. The record name \code{-} writes
the header to the standard output. A record name that begins with
\code{+} is not accepted (the Programmer's Guide says that the \code{+} is
removed; the code rejects the character).
\begin{params}
\item[record] The record name.
\end{params}
\returns{0 on success; -1 if the record name is illegal or the header cannot
be created.}
\raises{\code{init: illegal character N in record name}; \code{newheader:
can't create header for record REC}.}

\apientry{setheader}{Write a header from signal specifications}
\begin{lstlisting}[language=C,style=signature]
int setheader(char *record, const WFDB_Siginfo *siarray, unsigned int nsig);
\end{lstlisting}
Writes (or rewrites) the header of \code{record} for the \code{nsig} signals
described by \code{siarray}, with the frequencies, base time, skews and prolog
lengths as for \code{newheader}; the length is the \code{nsamp} of the first
signal. It serves to change a header, for example the gain of a signal: the
specifications are read with \code{isigopen(record, siarray, -nsig)}, changed
and written back; the info strings must be copied explicitly
(\code{getinfo}, \code{putinfo}). As in C, and although the argument is
declared \code{const}, the \code{units} of each structure are cut at their
first blank and the \code{desc} at their first line break, in the caller's
structures (in R and MATLAB/Octave, the changed structures are returned:
\code{list(stat, si)}, \code{[stat, si]}).
\begin{params}
\item[record] The record name.
\item[siarray] The signal specifications.
\item[nsig] The number of signals.
\end{params}
\returns{0 on success; -1 if the record name is illegal or the header cannot
be created.}
\raises{\code{newheader: can't create header for record REC}.}

\apientry{setmsheader}{Write the header of a multi-segment record}
\begin{lstlisting}[language=C,style=signature]
int setmsheader(char *record, char **segnames, unsigned int nsegments);
\end{lstlisting}
Closes the input signals, and writes the header of the multi-segment record
\code{record}, made of the \code{nsegments} existing records named in
\code{segnames}, in that order. The segments' signal files are neither
copied nor changed; the readers of the record (\code{getvec},
\code{isigsettime}, ...) go from one segment to the next. The segments must
have the same number of signals and the same sampling frequency, and headers
that give their length; the record takes the counter frequency, base counter
value and base time and date of the first segment.
\begin{params}
\item[record] The name of the multi-segment record.
\item[segnames] The names of the segments.
\item[nsegments] The number of segments.
\end{params}
\returns{0 on success; -1 if the record name is illegal, there is no segment,
or the header cannot be created; -2 if a segment name is longer than
\code{WFDB\_MAXRNL} (50) characters; -3 if the header of a segment cannot be
read (or is itself that of a multi-segment record); -4 if a segment's length
is not given, or its number of signals or sampling frequency differs from
those of the first segment.}
\raises{\code{setmsheader: record must contain at least one segment};
\code{setmsheader: can't read segment SEG header}; \code{setmsheader: length
of segment SEG must be specified}; \code{setmsheader: incorrect number of
signals in segment SEG}; \code{setmsheader: incorrect sampling frequency in
segment SEG}; \code{setmsheader: can't create header file for record REC}.}

\apientry{getseginfo}{The segments of the open record}
\begin{lstlisting}[language=C,style=signature]
int getseginfo(WFDB_Seginfo **segments);
\end{lstlisting}
Makes \code{*segments} point to the array of \code{WFDB\_Seginfo} structures
(\cref{sec:library-structures}) that describe the segments of the open
multi-segment record: the name, the length and the first sample number of
each. A segment named \code{\textasciitilde} is a gap, in which no signal is
available. If no multi-segment record is open, \code{*segments} is not
changed.
\begin{params}
\item[segments] The address of a pointer, set by the function.
\end{params}
\returns{The number of segments; 0 if the open record is not a
multi-segment record.}

The form differs between languages: in C++, \code{getseginfo(const
WFDB\_Seginfo **)} gives a pointer to the library's table; in Rust, Julia and
Kotlin, the argument (\code{\&mut Option<Vec<Seginfo>>}, a \code{Ref}, a
\code{MutableList<Seginfo>?}) receives a copy, and the number of segments is
returned; in Python, R and MATLAB/Octave, \code{getseginfo()} takes no
argument and returns the list (a struct array in MATLAB/Octave) of segments,
empty if the record is not a multi-segment record.

\apientry{getinfo}{Read the info strings of a record}
\begin{lstlisting}[language=C,style=signature]
char *getinfo(char *record);
\end{lstlisting}
With a record name, reads the info strings of \code{record} (the comment
lines of its header after the record line, then the lines of its
\file{.info} file that begin with \code{\#}, if it has one) and returns the
first; with \code{NULL}, returns the next one, reading those of the current
record (the last one opened or named) on the first call. An info string is
the line without its \code{\#} and its line ending.
\begin{params}
\item[record] The record name, or \code{NULL}.
\end{params}
\returns{The info string; \code{NULL} if there are no more.}
\raises{\code{getinfo: caller did not specify record name}.}

WFDB-x: the header and the \file{.info} file are read only where the record
name says (\cref{chap:behavior}).

\apientry{putinfo}{Write an info string}
\begin{lstlisting}[language=C,style=signature]
int putinfo(const char *info);
\end{lstlisting}
Writes \code{info} as an info string (a line \code{\#}\emph{info}, with a
CR-LF line ending): to the \file{.info} file opened by \code{setinfo}, or
else to the header being written (after \code{newheader}, \code{setheader}
or \code{setmsheader}). It should not contain a line break; WFDB-x does not
limit its length (the Programmer's Guide gives 254 characters).
\begin{params}
\item[info] The info string.
\end{params}
\returns{0 on success; -1 if no info file or header is open.}
\raises{\code{putinfo: caller has not specified a record name}.}

\apientry{setinfo}{Open the info file of a record}
\begin{lstlisting}[language=C,style=signature]
int setinfo(char *record);
\end{lstlisting}
Closes the info file opened by an earlier \code{setinfo}, if any, and opens
the \file{.info} file of \code{record} for writing with \code{putinfo},
appending to it if it exists. Such a file holds information about the record
without changing its header. \code{setinfo(NULL)} only closes the info file.
\begin{params}
\item[record] The record name, or \code{NULL}.
\end{params}
\returns{0 on success; -1 if the record name is illegal or the file cannot be
opened; with \code{NULL}, -2 if the info file could not be written
completely.}
\raises{\code{setinfo: can't create info file for record REC};
\code{setinfo: write error in info file}.}

\apientry{wfdb\_freeinfo}{Release the info strings}
\begin{lstlisting}[language=C,style=signature]
void wfdb_freeinfo(void);
\end{lstlisting}
Releases the info strings read by \code{getinfo}: the next \code{getinfo}
reads them again, from the first. \code{wfdbquit} calls it. (The Programmer's
Guide says that it also closes the file opened by \code{setinfo}; the code
does not, and neither does WFDB-x: \code{setinfo(NULL)} and \code{wfdbquit}
close it.)

\apientry{wfdbgetskew}{Skew of an input signal}
\begin{lstlisting}[language=C,style=signature]
int wfdbgetskew(WFDB_Signal s);
\end{lstlisting}
Returns the skew of input signal \code{s} recorded in the header, in frames
(the header gives it in samples). \code{getvec} and \code{getframe} already
correct the skew; the function serves programs that rewrite a header and
must keep the skews (with \code{wfdbsetskew}).
\begin{params}
\item[s] The input signal number.
\end{params}
\returns{The skew in frames; 0 if \code{s} is not an input signal.}

\apientry{wfdbsetiskew}{Change the skew of an input signal}
\begin{lstlisting}[language=C,style=signature]
void wfdbsetiskew(WFDB_Signal s, int skew);
\end{lstlisting}
Changes the skew, in frames, with which input signal \code{s} is read (a
function of the C library that the Programmer's Guide does not document). The
new skew is applied only if it is between 0 and the largest skew of the
record's header; otherwise nothing changes.
\begin{params}
\item[s] The input signal number.
\item[skew] The skew, in frames.
\end{params}

\apientry{wfdbsetskew}{Set the skew of an output signal}
\begin{lstlisting}[language=C,style=signature]
void wfdbsetskew(WFDB_Signal s, int skew);
\end{lstlisting}
Sets the skew, in frames, that \code{newheader} and \code{setheader} write
for output signal \code{s} (in samples). It does not affect the reading of
input signals. Nothing is done if \code{s} is not an output signal.
\begin{params}
\item[s] The output signal number.
\item[skew] The skew, in frames.
\end{params}

\apientry{wfdbgetstart}{Prolog length of an input signal file}
\begin{lstlisting}[language=C,style=signature]
long wfdbgetstart(WFDB_Signal s);
\end{lstlisting}
Returns the length of the prolog (the bytes before sample 0) of the signal
file of input signal \code{s}, as the header records it; the readers already
skip it. With no input signals open, signal 0 gives the prolog of the first
signal of the last header read.
\begin{params}
\item[s] The input signal number.
\end{params}
\returns{The prolog length in bytes; 0 if \code{s} is not an input signal.}

\apientry{wfdbsetstart}{Set the prolog length of an output signal}
\begin{lstlisting}[language=C,style=signature]
void wfdbsetstart(WFDB_Signal s, long bytes);
\end{lstlisting}
Sets the prolog length that \code{newheader} and \code{setheader} write for
the signal file of output signal \code{s}; it does not affect the reading of
input signals. The value also applies to every signal of the next header
written by \code{setheader} whose file has no prolog length of its own (as
when no output signals are open).
\begin{params}
\item[s] The output signal number.
\item[bytes] The prolog length, in bytes.
\end{params}

\apientry{wfdbputprolog}{Write a prolog to an output signal file}
\begin{lstlisting}[language=C,style=signature]
int wfdbputprolog(const char *prolog, long bytes, WFDB_Signal s);
\end{lstlisting}
Writes the first \code{bytes} bytes of \code{prolog} to the signal file of
output signal \code{s} (normally before any sample), and records the number
of bytes written with \code{wfdbsetstart}.
\begin{params}
\item[prolog] The bytes of the prolog.
\item[bytes] Their number.
\item[s] The output signal number.
\end{params}
\returns{0 on success; -1 if fewer bytes were written.}
\raises{\code{wfdbputprolog: only N of M bytes written}.}

% ==========================================================================
\section{Error handling and settings}
\label{sec:library-errors}

The library reports an error by its return value and by a message, printed on
the standard error output unless \code{wfdbquiet} was called; the last
message is kept for \code{wfdberror}.

\apientry{wfdbquiet}{Stop printing error messages}
\begin{lstlisting}[language=C,style=signature]
void wfdbquiet(void);
\end{lstlisting}
Stops the printing of the library's messages on the standard error output;
they are still kept for \code{wfdberror}.

\apientry{wfdbverbose}{Print error messages again}
\begin{lstlisting}[language=C,style=signature]
void wfdbverbose(void);
\end{lstlisting}
Restores the printing of the library's messages (the default).

\apientry{wfdberror}{The last error message}
\begin{lstlisting}[language=C,style=signature]
char *wfdberror(void);
\end{lstlisting}
Returns the last message of the library, printed or not.
\returns{The message, with its final newline; if there has been none, the
library version, \code{"WFDB library version 10.7.0 (Jun 10 2022).\textbackslash
n"} (the build date of the C release reproduced).}

\apientry{wfdb\_error}{Report an error as the library does}
\begin{lstlisting}[language=C,style=signature]
void wfdb_error(const char *format_string, ...);
\end{lstlisting}
Formats a message as \code{printf} does, keeps it for \code{wfdberror} and
prints it on the standard error output unless \code{wfdbquiet} was called.
The library reports its own errors with it; a program may use it for its
own messages.
\begin{params}
\item[format\_string] A \code{printf} format.
\item[...] The values to format.
\end{params}

In Python, \code{wfdb\_error(fmt, *args)}; in Rust and Kotlin,
\code{wfdb\_error(msg)} takes a message already formatted (Rust also has
\code{wfdb\_error\_fmt(fmt, args)}); in Julia, \code{WFDBX.wfdb\_error(fmt,
args...)} (not exported). The formatting is that of the GNU C library in
every language.

\apientry{wfdbmemerr}{Behavior on memory errors}
\begin{lstlisting}[language=C,style=signature]
void wfdbmemerr(int exit_on_error);
\end{lstlisting}
In C, chooses whether a failure to allocate memory ends the program (nonzero,
the default) or lets the function continue (0). WFDB-x keeps the value, for
\code{wfdb\_me\_fatal}, but has no other use for it: the languages report
memory exhaustion themselves.
\begin{params}
\item[exit\_on\_error] Nonzero to end the program on a memory error.
\end{params}

\apientry{wfdb\_me\_fatal}{Whether memory errors are fatal}
\begin{lstlisting}[language=C,style=signature]
int wfdb_me_fatal(void);
\end{lstlisting}
Returns the value set by \code{wfdbmemerr}.
\returns{The value; 1 by default.}

\apientry{wfdbversion}{Version of the library}
\begin{lstlisting}[language=C,style=signature]
const char *wfdbversion(void);
\end{lstlisting}
Returns the version of the WFDB library that WFDB-x reproduces (the version
of WFDB-x itself is that of its packages, \wfdbxversion).
\returns{\code{"10.7.0"}.}

\apientry{wfdbldflags}{Linker flags of the C library}
\begin{lstlisting}[language=C,style=signature]
const char *wfdbldflags(void);
\end{lstlisting}
In C, returns the linker options with which programs are linked with the
library (fixed when it is built). WFDB-x has none.
\returns{\code{""}.}

\apientry{wfdbcflags}{Compiler flags of the C library}
\begin{lstlisting}[language=C,style=signature]
const char *wfdbcflags(void);
\end{lstlisting}
In C, returns the compiler options with which programs that use the library
are compiled. WFDB-x has none.
\returns{\code{""}.}

\apientry{wfdbdefwfdb}{Default WFDB path}
\begin{lstlisting}[language=C,style=signature]
const char *wfdbdefwfdb(void);
\end{lstlisting}
Returns the default WFDB path of the C library, used in legacy mode when the
environment variable \code{WFDB} is not set (that of a C release installed in
\file{/usr/local}).
\returns{\code{". /usr/local/database http://physionet.org/physiobank/database"}.}

\apientry{wfdbdefwfdbcal}{Default calibration file}
\begin{lstlisting}[language=C,style=signature]
const char *wfdbdefwfdbcal(void);
\end{lstlisting}
Returns the name of the calibration file that \code{calopen(NULL)} reads when
the environment variable \code{WFDBCAL} is not set.
\returns{\code{"wfdbcal"}.}

% ==========================================================================
\section{File and path functions}
\label{sec:library-files}

The C library finds its input files by searching the WFDB path, a list of
directories and URLs. WFDB-x finds a file only where its name says, under the
input directory for a relative name (\cref{chap:behavior,chap:files}): the
functions of this section keep their C names and arguments, and set and
return the input directory; in legacy mode, they work on the WFDB path as in
C, without network access.

\apientry{setwfdb}{Set the input directory (the WFDB path)}
\begin{lstlisting}[language=C,style=signature]
void setwfdb(const char *database_path_string);
\end{lstlisting}
WFDB-x: sets the input directory. The string is split as a WFDB path is (at
spaces, tabs, newlines, \code{;}, and \code{:} outside a URL), and must have at
most one nonempty component, which becomes the input directory (a directory,
a \code{\textasciitilde} path or a URL); with none (\code{""}, \code{";"},
\code{NULL}), the default is restored (\code{WFDBX\_INPUT\_DIR}, or the
current directory). With more than one, nothing changes and an error is
reported. A directory whose name contains a space must be set with
\code{wfdbx\_set\_input\_dir}.

In legacy mode, as in C: sets the WFDB path, the list of locations searched
for input files (\code{@file} reads it from a file, up to ten levels deep);
\code{NULL} restores the value of the environment variable \code{WFDB}, or
the default (\code{wfdbdefwfdb}); the value is also exported as the
environment variable \code{WFDB}, for the programs started afterwards; URL
components are not used.
\begin{params}
\item[database\_path\_string] The input directory (legacy mode: the WFDB
  path), or \code{NULL}.
\end{params}
\raises{\code{setwfdb: WFDB-x reads from a single input directory; use
wfdbx\_set\_input\_dir (or legacy mode for the WFDB search path)}.}

\apientry{getwfdb}{The input directory (the WFDB path)}
\begin{lstlisting}[language=C,style=signature]
char *getwfdb(void);
\end{lstlisting}
WFDB-x: returns the input directory. In legacy mode, as in C: returns the
WFDB path, read on the first call from the environment variable \code{WFDB}
(or the default), \code{@file} indirections evaluated. The environment
variable \code{WFDB} is ignored outside legacy mode, and the library says so
once (\code{WFDB-x: the WFDB environment variable is ignored; set
WFDBX\_INPUT\_DIR instead (or use legacy mode)}).
\returns{The input directory, or \code{"."} if none is set (legacy mode: the
WFDB path).}

\apientry{resetwfdb}{Restore the input directory (the WFDB path)}
\begin{lstlisting}[language=C,style=signature]
void resetwfdb(void);
\end{lstlisting}
WFDB-x: restores the input directory set by the last \code{setwfdb}, if any
(\code{wfdbx\_set\_input\_dir} may have changed it since). In legacy mode, as
in C: restores the WFDB path to the value that the first \code{getwfdb} of the
process returned, if \code{getwfdb} was called before \code{setwfdb};
\code{wfdbquit} calls it (in legacy mode only, in WFDB-x).

\apientry{wfdbfile}{The path of a WFDB file}
\begin{lstlisting}[language=C,style=signature]
char *wfdbfile(const char *file_type, char *record);
\end{lstlisting}
Returns the name under which the library opens the file of \code{record}
with suffix \code{file\_type} (\code{"hea"}, \code{"atr"}, ...), if it can be
opened: \code{wfdbfile("hea", "100s")} gives \file{100s.hea}, or the path
under the input directory. A file whose name does not contain a record name
is given by \code{file\_type} with \code{record} \code{NULL}. As in C, a
second name, with the suffix cut to 3 characters, is tried. With both
arguments \code{NULL}, it returns the name of the file most recently opened
by the library.
\begin{params}
\item[file\_type] The file type (suffix), or a file name.
\item[record] The record name, or \code{NULL}.
\end{params}
\returns{The path name or URL; \code{NULL} if the file cannot be opened.}

WFDB-x: the name is the one location that the name designates (an absolute
name or URL as given, a relative name under the input directory), not the
first found in the WFDB path; a URL is opened only with network access
(\cref{chap:behavior}).

% ==========================================================================
\section{WFDB-x additions}
\label{sec:library-additions}

These functions are not in the C library. Their names begin with
\code{wfdbx\_}, except the reader of many sample vectors; they are the same
in the seven languages, and the signatures below are those of the C++
implementation (\file{cpp/include/wfdbx/wfdb.hpp}). \Cref{chap:behavior}
explains the settings they control; the environment variables
\code{WFDBX\_INPUT\_DIR}, \code{WFDBX\_OUTPUT\_DIR} and \code{WFDBX\_NETWORK}
give their initial values, which a call overrides.

\apientry{wfdbx\_set\_input\_dir}{Set the input directory}
\begin{lstlisting}[language=C++,style=signature]
std::optional<std::string> wfdbx_set_input_dir(const char *dir);
\end{lstlisting}
Sets the input directory: every file named by a relative name (\code{100s},
\code{mitdb/100}, \code{../other/100}) is read from it; an absolute name, a
\code{\textasciitilde} name or a URL is used as given. It may be a local
directory, a \code{\textasciitilde} path, or the URL of a directory, such as
\code{https://physionet.org/files/mitdb/1.0.0/} (which also needs network
access, \code{wfdbx\_set\_network}). The name is used as it is (it may contain
spaces). The setting applies to the library functions, the high-level
functions (whose \code{input\_dir} option sets it for the duration of a call)
and the tools, also for the files that they read themselves
(\code{wfdbx\_input\_path}); \code{wfdbquit} does not change it; legacy mode
does not use it (it searches the WFDB path).
\begin{params}
\item[dir] The input directory; \code{NULL} or \code{""} restores the
  default: the value of \code{WFDBX\_INPUT\_DIR} if it is set, otherwise the
  current directory.
\end{params}
\returns{The previous input directory (that of \code{WFDBX\_INPUT\_DIR} if
it was in effect); \code{NULL} (\code{std::nullopt}, \code{None},
\code{nothing}, ...) if none was set.}

\apientry{wfdbx\_get\_input\_dir}{The input directory}
\begin{lstlisting}[language=C++,style=signature]
std::optional<std::string> wfdbx_get_input_dir();
\end{lstlisting}
Returns the input directory set by \code{wfdbx\_set\_input\_dir},
\code{setwfdb} or \code{WFDBX\_INPUT\_DIR}.
\returns{The input directory; \code{NULL} if none is set (files are then read
from the current directory; \code{getwfdb} returns \code{"."}).}

\apientry{wfdbx\_set\_output\_dir}{Set the output directory}
\begin{lstlisting}[language=C++,style=signature]
std::optional<std::string> wfdbx_set_output_dir(const char *dir);
\end{lstlisting}
Sets the output directory: every file written under a relative name (a
signal file, a header, an annotation, info or calibration file) is created
in it; missing directories are created. An absolute name or a
\code{\textasciitilde} name is used as given. The names recorded in headers
do not change. The setting applies to the library functions, the high-level
writers (whose \code{dir} option sets it for the duration of a call) and the
tools, also for the files that they create themselves
(\code{wfdbx\_output\_path}); \code{wfdbquit} does not change it. In legacy
mode it is a plain prefix of every output name, absolute names included.
\begin{params}
\item[dir] A local directory or a \code{\textasciitilde} path; \code{NULL} or
  \code{""} restores the default: the value of \code{WFDBX\_OUTPUT\_DIR} if it
  is set, otherwise the current directory. A URL is refused.
\end{params}
\returns{The previous output directory; \code{NULL} if none was set. For a
URL, the current output directory, unchanged.}
\raises{\code{wfdbx\_set\_output\_dir: the output directory must be a local
directory}.}

\apientry{wfdbx\_get\_output\_dir}{The output directory}
\begin{lstlisting}[language=C++,style=signature]
std::optional<std::string> wfdbx_get_output_dir();
\end{lstlisting}
Returns the output directory set by \code{wfdbx\_set\_output\_dir} or
\code{WFDBX\_OUTPUT\_DIR}.
\returns{The output directory; \code{NULL} if none is set (files are then
written in the current directory).}

\apientry{wfdbx\_set\_network}{Allow or forbid network access}
\begin{lstlisting}[language=C++,style=signature]
int wfdbx_set_network(int on);
\end{lstlisting}
Allows or forbids reading files from \code{http://} and \code{https://}
URLs; it is forbidden by default. A remote file is read in full the first
time it is opened and kept for the rest of the process; a download gives up
if it cannot connect, or receives nothing, for 30 seconds. With network
access forbidden, the first attempt to open a URL prints \code{WFDB-x:
network access is disabled; set WFDBX\_NETWORK=1 or call wfdbx\_set\_network(1)
to open URL}, and the file cannot be opened. Legacy mode reads no remote
files.
\begin{params}
\item[on] 1 (or any positive value; \code{True} in Python) to allow, 0 to
  forbid; a negative value (or \code{None} in Python) restores the default:
  allowed if \code{WFDBX\_NETWORK} is a nonzero number.
\end{params}
\returns{The previous setting: 1 (allowed) or 0.}

\apientry{wfdbx\_get\_network}{Whether network access is allowed}
\begin{lstlisting}[language=C++,style=signature]
int wfdbx_get_network();
\end{lstlisting}
\returns{1 if network access is allowed, 0 otherwise.}

\apientry{wfdbx\_set\_legacy}{Turn legacy mode on or off}
\begin{lstlisting}[language=C++,style=signature]
int wfdbx_set_legacy(int on);
\end{lstlisting}
Turns legacy mode on or off for the whole library. In legacy mode, every
deliberate change of WFDB-x is undone (\cref{chap:behavior}): the built-in
sampling frequency, gain and ADC resolution, the WFDB search path (without
network access), sorting by the \code{sortann} program, output files that may
be input files, and C's checks and names in headers. It is off by default,
and is meant only to reproduce results of the original package.
\begin{params}
\item[on] Nonzero to turn legacy mode on, 0 to turn it off.
\end{params}
\returns{The previous setting: 1 (on) or 0.}

\apientry{wfdbx\_get\_legacy}{Whether legacy mode is on}
\begin{lstlisting}[language=C++,style=signature]
int wfdbx_get_legacy();
\end{lstlisting}
\returns{1 if legacy mode is on, 0 otherwise.}

\apientry{wfdbx\_legacy\_args}{Recognize the --legacy option of a tool}
\begin{lstlisting}[language=C++,style=signature]
int wfdbx_legacy_args(int &argc, char **&argv);
\end{lstlisting}
The function that every tool calls first: if the first argument after the
program name (\code{argv[1]}) is exactly \code{--legacy}, it is removed (the
program name stays first) and legacy mode is turned on. \code{--legacy}
elsewhere is not recognized. In C++, Rust and Kotlin, the call also arranges
the exit status of the program as for the tools: at exit, outside legacy
mode, a failure to write the standard output is reported (\code{PROGRAM:
write error on standard output}), and an exit status of 0 becomes 1 after a
write error (see \code{wfdbx\_write\_errors}).
\begin{params}
\item[argc] The number of arguments, decremented if \code{--legacy} is
  removed.
\item[argv] The arguments, program name first.
\end{params}
\returns{1 if \code{--legacy} was found (and removed), 0 otherwise.}

The argument list is that of each language: in Python, a list
(\code{sys.argv}), changed in place; in Rust and Kotlin, a mutable list of
byte strings; in R, a character vector, and the result is \code{list(argv,
legacy)}; in MATLAB/Octave, a cell array, with \code{[r, argv] =
wfdbx.wfdbx\_legacy\_args(argv)}. In Julia, \code{wfdbx\_legacy\_args(args)}
takes the arguments \emph{without} the program name (Julia's \code{ARGS}), and
looks at \code{args[1]}.

\apientry{wfdbx\_time\_needs\_freq}{Whether strtim needs the sampling frequency}
\begin{lstlisting}[language=C++,style=signature]
int wfdbx_time_needs_freq(const char *time_string);
\end{lstlisting}
Tells whether \code{strtim} needs a sampling frequency to convert
\code{time\_string}, so that a program can report a missing frequency in its
own terms.
\begin{params}
\item[time\_string] The time string.
\end{params}
\returns{0 if the string (after leading spaces, tabs, carriage returns and
newlines) begins with \code{e}, \code{i}, \code{o} or \code{s}, or with
\code{c} while no counter frequency is defined; 1 otherwise.}

\apientry{wfdbx\_write\_errors}{Number of write errors}
\begin{lstlisting}[language=C++,style=signature]
int wfdbx_write_errors();
\end{lstlisting}
Returns the number of write errors since the program started: the output
files (signal files, headers, annotation, info and calibration files, and
the standard output opened as \code{-}) that the library could not write,
flush or close, found when they are closed (by \code{wfdbquit},
\code{osigfopen(NULL, 0)}, \code{oannclose}, ...), and those counted by
\code{wfdbx\_note\_write\_error}. The count belongs to the process, not to a
library state object (Rust, Julia, Kotlin). The tools exit with status 1
when it is not 0, and the high-level writers raise an error when they meet a
write error (\cref{chap:behavior}).
\returns{The number of write errors.}

\apientry{wfdbx\_note\_write\_error}{Count a write error}
\begin{lstlisting}[language=C++,style=signature]
void wfdbx_note_write_error();
\end{lstlisting}
Adds one to the count of \code{wfdbx\_write\_errors}, for a program that
writes a file itself and could not write it. (Rust and Kotlin also have
\code{wfdbx\_fclose\_output(file, pname, name)}, which closes such a file,
and reports \code{PNAME: write error on NAME} and counts it on failure.)

\apientry{wfdbx\_input\_path}{Where a program reads a file itself}
\begin{lstlisting}[language=C++,style=signature]
std::optional<std::string> wfdbx_input_path(const char *name);
\end{lstlisting}
Returns the path under which a program opens a file that it reads itself,
rather than through the library, so that the file is found where
\code{wfdb\_open} finds a file of that name (with an empty file type). In
normal mode, a relative name is found under the input directory, if one is
set (\code{wfdbx\_set\_input\_dir}), a \code{\textasciitilde} name is
expanded, and an absolute name is used as given; in legacy mode, the name is
used as given, as the original programs did. The tools use it for the files
named on their command line that they read themselves, such as the \opt{-i}
file of \cmd{wrsamp} (\cref{sec:tool-files}). For example, with the input
directory \file{data}, \code{wfdbx\_input\_path("samples.txt")} returns
\code{data/samples.txt}.
\begin{params}
\item[name] A file name. \code{""} and \code{"-"} (the standard input, if the
  program treats it so) are returned as given.
\end{params}
\returns{The path to open; \code{NULL} (\code{std::nullopt}, \code{None},
\code{nothing}, ...) if \code{name} is \code{NULL} or a URL, or becomes a URL
when the input directory is prefixed to it: such a file cannot be opened in
this way.}

\apientry{wfdbx\_output\_path}{Where a program creates a file itself}
\begin{lstlisting}[language=C++,style=signature]
std::optional<std::string> wfdbx_output_path(const char *name);
\end{lstlisting}
Returns the path under which a program creates a file that it writes
itself, rather than through the library, so that the file is created where
\code{wfdb\_open} creates a file of that name (with an empty file type). In
normal mode, a relative name is created under the output directory, if one
is set (\code{wfdbx\_set\_output\_dir}), a \code{\textasciitilde} name is
expanded, and an absolute name is used as given; in legacy mode, the output
directory, if one is set, is a plain prefix of any name, absolute names
included. When the output directory has been prefixed to the name, the
directories missing in the resulting path are created, as for the output
files of the library. The tools use it for the files that they create
themselves, such as the WAV file of \cmd{mit2wav} (\cref{sec:tool-files}).
For example, with the output directory \file{out},
\code{wfdbx\_output\_path("notes.txt")} returns \code{out/notes.txt}.
\begin{params}
\item[name] A file name. \code{""} and \code{"-"} (the standard output, if the
  program treats it so) are returned as given.
\end{params}
\returns{The path to create; \code{NULL} (\code{std::nullopt}, \code{None},
\code{nothing}, ...) if \code{name} is \code{NULL} or a URL (which cannot be
written).}

\apientry{getframes}{Read many sample vectors at once}
\begin{lstlisting}[language=C++,style=signature]
long long getvecs(WFDB_Sample *out, long long n, int *status = nullptr);
// The same function in each language:
//   C++, Rust, Julia, Kotlin:  getvecs(out, n), and in Rust, Julia and Kotlin
//                              getvecs_status(out, n), which also returns status
//   Python:                    frames, status = getframes(n)
//   R:                         r <- getframes(n, minstat = 0)   # r$v, r$stat
//   MATLAB/Octave:             [v, status, last] = wfdbx.getframes(n, minstat)
\end{lstlisting}
Reads up to \code{n} sample vectors, with exactly the result of \code{n}
calls of \code{getvec} that stop at the first negative status, but much
faster: whole blocks of the signal files are decoded at once, for every
format, for multi-frequency records in both modes of \code{getvec}, and for
multi-segment records; the frames that need special handling (the one at
which a checksum is verified, a change of segment, skewed or resampled
signals) go through \code{getvec}. The vectors follow one another in
\code{out}, one value per signal in each (in Python, R and MATLAB/Octave, a
matrix with one row per vector and one column per signal).
\begin{params}
\item[out] (C++, Rust, Julia, Kotlin) An array with room for \code{n} times
  the number of values of a vector, filled by the function.
\item[n] The maximum number of vectors to read (\code{Inf} in R and
  MATLAB/Octave reads to the end).
\item[status] (C++, optional) Receives the value returned by the last
  \code{getvec}: the negative status of the first vector that could not be
  read (-1 at the end of the data, -3 for an unexpected end of file, ...), or
  the number of signals if all \code{n} were read.
\item[minstat] (R, MATLAB/Octave, optional) Stop at the first vector for
  which \code{getvec} would return less than \code{minstat} (as a C loop
  \code{while (getvec(v) >= minstat)}); 0 by default.
\end{params}
\returns{C++, Rust, Julia, Kotlin: the number of vectors read (\code{>= 0}),
or the negative status of \code{getvec} if none could be read;
\code{getvecs\_status} returns the number and the status. Python: a NumPy
array (\code{int64}, one row per vector) and the status. R:
\code{list(v, stat)}, a matrix and the status. MATLAB/Octave: the matrix, the
status, and the vector filled by the \code{getvec} call that stopped the
reading (the last frame of the record for a checksum error, -4).}

% ==========================================================================
\section{Constants and structures}
\label{sec:library-structures}

The constants of \file{<wfdb/wfdb.h>} have their C names and values in every
language (\cref{tab:library-constants}); in MATLAB/Octave they are functions
without arguments (\code{wfdbx.WFDB\_READ()}), and the calibration constants
and \code{WFDB\_FMT\_LIST} are not defined there. The annotation codes
(\code{NORMAL}, \code{PVC}, ..., \code{ACMAX}) of \file{<wfdb/ecgcodes.h>}
are listed in \cref{chap:formats}.

\begin{xltabular}{\linewidth}{@{}P{4.4cm}P{2.4cm}L@{}}
\caption{The constants of the library.}
\label{tab:library-constants}\\
\toprule
{\tabhead Constant} & {\tabhead Value} & {\tabhead Meaning}\\
\midrule
\endfirsthead
\toprule
{\tabhead Constant} & {\tabhead Value} & {\tabhead Meaning}\\
\midrule
\endhead
\bottomrule
\endfoot
\code{WFDB\_READ} & 0 & \code{stat} of an MIT-format input annotation file\\
\code{WFDB\_WRITE} & 1 & \code{stat} of an MIT-format output annotation file\\
\code{WFDB\_AHA\_READ} & 2 & \code{stat} of an AHA-format input annotation
file\\
\code{WFDB\_AHA\_WRITE} & 3 & \code{stat} of an AHA-format output annotation
file\\
\code{WFDB\_APPEND} & 4 & mode of an output info file (opened for appending)\\
\code{WFDB\_INVALID\_SAMPLE} & -32768 & value of an invalid or missing
sample, in \code{getvec}, \code{getframe} and \code{sample}, unless
\code{WFDB\_GVPAD} is set; written by \code{putvec} as the invalid value of
the format\\
\code{WFDB\_LOWRES} & 0 & \code{getvec} mode: one sample per signal per
frame\\
\code{WFDB\_HIGHRES} & 1 & \code{getvec} mode: every sample of the
oversampled signals\\
\code{WFDB\_GVPAD} & 2 & \code{getvec} flag: invalid samples replaced by the
previous valid sample\\
\code{WFDB\_DEFFREQ} & 250.0 & C's default sampling frequency (Hz); not used
by WFDB-x outside legacy mode\\
\code{WFDB\_DEFGAIN} & 200.0 & C's default gain (adu per physical unit); not
used outside legacy mode\\
\code{WFDB\_DEFRES} & 12 & C's default ADC resolution (bits); not used
outside legacy mode\\
\code{WFDB\_MAXRNL} & 50 & maximum length of a record name (of a segment
name); not exported in R\\
\code{WFDB\_MAXUSL} & 50 & maximum length of a \code{units} string\\
\code{WFDB\_MAXDSL} & 100 & maximum length of a \code{desc} string\\
\code{WFDB\_FMT\_LIST} & 0, 8, 16, 61, 80, 160, 212, 310, 311, 24, 32, 508,
516, 524 & the signal storage formats (\cref{chap:formats})\\
\code{WFDB\_AC\_COUPLED} & 0 & \code{caltype}: AC-coupled signal\\
\code{WFDB\_DC\_COUPLED} & 1 & \code{caltype}: DC-coupled signal\\
\code{WFDB\_CAL\_SQUARE} & 2 & \code{caltype}: square calibration pulse\\
\code{WFDB\_CAL\_SINE} & 4 & \code{caltype}: sine calibration pulse\\
\code{WFDB\_CAL\_SAWTOOTH} & 6 & \code{caltype}: sawtooth calibration pulse\\
\code{WFDB\_CAL\_UNDEF} & 8 & \code{caltype}: undefined pulse shape\\
\code{APUNDEF}, \code{APSTD}, \code{APHIGH}, \code{APLOW}, \code{APATT},
\code{APAHIGH}, \code{APALOW} & 0 to 6 & display positions (\code{annpos}):
undefined, standard, above, below, attached to the signal \code{chan}, above
and below that\\
\end{xltabular}

The simple types of the C library are integers and floating-point numbers:
\code{WFDB\_Sample} (\code{int}, 32 bits) for samples in ADC units,
\code{WFDB\_Time} (64 bits on the reference platform) for times in sample
intervals, \code{WFDB\_Date} (\code{long}) for Julian dates,
\code{WFDB\_Frequency} and \code{WFDB\_Gain} (\code{double}), and
\code{WFDB\_Group}, \code{WFDB\_Signal} and \code{WFDB\_Annotator}
(\code{unsigned int}). \Cref{sec:library-calling} gives the corresponding
types of each language (doubles in R and MATLAB/Octave).

The structures have the C field names in every language
(\cref{tab:library-siginfo,tab:library-anninfo,tab:library-annotation,tab:library-calinfo,tab:library-seginfo});
their names are \code{WFDB\_Siginfo}, ..., in C++ (also \code{Siginfo}, ...),
\code{Siginfo}, \code{Anninfo}, \code{Annotation}, \code{Calinfo} and
\code{Seginfo} in Python, Rust, Julia and Kotlin, and lists or structures made
by \code{siginfo()}, \code{anninfo()}, \code{annotation()}, \code{calinfo()}
and \code{seginfo()} in R and MATLAB/Octave (\code{wfdbx.siginfo(n)}, ...). A
string field that C may leave \code{NULL} is absent in the language's terms
(\code{None}, \code{std::nullopt}, \code{nothing}, \code{NULL}, \code{[]},
\code{null}).

\begin{xltabular}{\linewidth}{@{}P{3.4cm}L@{}}
\caption{\code{WFDB\_Siginfo}: the specification of a signal.}
\label{tab:library-siginfo}\\
\toprule
{\tabhead Field} & {\tabhead Meaning}\\
\midrule
\endfirsthead
\toprule
{\tabhead Field} & {\tabhead Meaning}\\
\midrule
\endhead
\bottomrule
\endfoot
\code{char *fname} & the name of the signal file, relative to the header's
directory (WFDB-x); \code{-} is the standard input or output\\
\code{char *desc} & the description of the signal (\code{MLII}), at most
\code{WFDB\_MAXDSL} characters, without line breaks; a header that gives none
gives \code{record REC, signal N}\\
\code{char *units} & the physical units, at most \code{WFDB\_MAXUSL}
characters, without blanks; \code{NULL} if not specified (WFDB-x: never
taken to be millivolts)\\
\code{WFDB\_Gain gain} & ADC units per physical unit; 0 if uncalibrated\\
\code{WFDB\_Sample initval} & the value of sample 0 (used by format 8)\\
\code{WFDB\_Group group} & the signal group: the signals of a group are stored
in one file, and have consecutive numbers\\
\code{int fmt} & the storage format (\code{WFDB\_FMT\_LIST},
\cref{chap:formats})\\
\code{int spf} & samples per frame (more than 1 for an oversampled signal)\\
\code{int bsize} & the block size in bytes, for character special files (0
for ordinary files)\\
\code{int adcres} & the ADC resolution in bits; 0 if not specified (WFDB-x:
not inferred from the format)\\
\code{int adczero} & the ADC output for an input at the center of its range\\
\code{int baseline} & the ADC value of 0 physical units\\
\code{long nsamp} & the number of samples per signal (of frames, for an
oversampled signal); 0 if not specified\\
\code{int cksum} & the 16-bit checksum of the samples, verified by
\code{getvec} at the end of the record\\
\end{xltabular}

The constructors give C's zero values, except that \code{spf} is 1 in Python,
R and MATLAB/Octave (0 in C++, Rust, Julia and Kotlin; \code{osigfopen} takes
an \code{spf} below 1 as 1).

\begin{xltabular}{\linewidth}{@{}P{3.4cm}L@{}}
\caption{\code{WFDB\_Anninfo}: an annotator to open.}
\label{tab:library-anninfo}\\
\toprule
{\tabhead Field} & {\tabhead Meaning}\\
\midrule
\endfirsthead
\toprule
{\tabhead Field} & {\tabhead Meaning}\\
\midrule
\endhead
\bottomrule
\endfoot
\code{char *name} & the annotator name, the suffix of the annotation file
(\code{atr} is the reference annotation file of a database record);
\code{-} is the standard input or output\\
\code{int stat} & \code{WFDB\_READ}, \code{WFDB\_WRITE},
\code{WFDB\_AHA\_READ} or \code{WFDB\_AHA\_WRITE}\\
\end{xltabular}

\begin{xltabular}{\linewidth}{@{}P{3.4cm}L@{}}
\caption{\code{WFDB\_Annotation}: an annotation.}
\label{tab:library-annotation}\\
\toprule
{\tabhead Field} & {\tabhead Meaning}\\
\midrule
\endfirsthead
\toprule
{\tabhead Field} & {\tabhead Meaning}\\
\midrule
\endhead
\bottomrule
\endfoot
\code{WFDB\_Time time} & the time, in sample intervals from the beginning of
the record (or in the units of \code{setiafreq} and \code{setafreq})\\
\code{char anntyp} & the annotation code, 1 to \code{ACMAX}
(\cref{chap:formats})\\
\code{signed char subtyp} & the subtype (-128 to 127); for noise
annotations, the signals affected\\
\code{unsigned char chan} & the channel (0 to 255), usually the signal to
which the annotation refers\\
\code{signed char num} & the annotator number (-128 to 127), free for the
program's use\\
\code{unsigned char *aux} & \code{NULL}, or the auxiliary data: its first
byte is the number of bytes that follow (at most 255), such as the rhythm of
a rhythm change \code{(N}, or the text of a comment; it may contain any
bytes\\
\end{xltabular}

\begin{xltabular}{\linewidth}{@{}P{3.4cm}L@{}}
\caption{\code{WFDB\_Calinfo}: an entry of the calibration list.}
\label{tab:library-calinfo}\\
\toprule
{\tabhead Field} & {\tabhead Meaning}\\
\midrule
\endfirsthead
\toprule
{\tabhead Field} & {\tabhead Meaning}\\
\midrule
\endhead
\bottomrule
\endfoot
\code{double low} & the low level of the calibration pulse, in physical
units (0 for an AC-coupled signal)\\
\code{double high} & the high level of the pulse (the pulse amplitude for an
AC-coupled signal)\\
\code{double scale} & the customary plotting scale, in physical units per
centimeter\\
\code{char *sigtype} & the signal type, matched against the beginning of a
signal description (empty: any signal)\\
\code{char *units} & the physical units\\
\code{int caltype} & \code{WFDB\_AC\_COUPLED} or \code{WFDB\_DC\_COUPLED},
plus \code{WFDB\_CAL\_SQUARE}, \code{WFDB\_CAL\_SINE} or
\code{WFDB\_CAL\_SAWTOOTH}\\
\end{xltabular}

\begin{xltabular}{\linewidth}{@{}P{3.4cm}L@{}}
\caption{\code{WFDB\_Seginfo}: a segment of a multi-segment record.}
\label{tab:library-seginfo}\\
\toprule
{\tabhead Field} & {\tabhead Meaning}\\
\midrule
\endfirsthead
\toprule
{\tabhead Field} & {\tabhead Meaning}\\
\midrule
\endhead
\bottomrule
\endfoot
\code{char recname[WFDB\_MAXRNL+1]} & the name of the segment (a record), or
\code{\textasciitilde} for a gap without signals\\
\code{WFDB\_Time nsamp} & the length of the segment, in samples\\
\code{WFDB\_Time samp0} & the number of samples before the segment in the
multi-segment record (the sample number of its first sample)\\
\end{xltabular}

\setlength{\emergencystretch}{0pt}
